% concatenated from RMCfromDM.tex, settings.tex, mymacros.tex
\documentclass[12pt]{amsart}

%Packages------------------------------------------------------

\usepackage[all]{xypic}
\usepackage{tikz}
\usepackage{ragged2e}
\usetikzlibrary{arrows} 
\usetikzlibrary{decorations.markings}
\usepackage{graphicx}
\usepackage{bm}
\usepackage{float}
\usepackage{epsf}
\usepackage{verbatim} 
\usepackage{amsmath}
\usepackage{amsfonts}
\usepackage{amssymb}
\usepackage{mathrsfs}
\usepackage{amsthm}
\usepackage{newlfont}
\usepackage{enumitem}
\usepackage[new]{old-arrows}
\usepackage{booktabs}
\usepackage{enumitem}
%\usepackage{amsmidx}
%\usepackage{makeidx} % this clashes with amsart etc since these have index built in
%\usepackage{stackengine}
%\usepackage{scalerel}
%\usepackage{amsaddr}
\usepackage{makecell}
\usepackage[multiple]{footmisc}
\usepackage[answerdelayed,lastexercise]{exercise}
\usepackage[hidelinks]{hyperref}
\usepackage{mleftright}

\usepackage{fnpct}
%\usepackage{fullpage} % From Poonen's "writing". This messes-up the book format
\usepackage{colonequals} % From Poonen's "writing"
%\usepackage{amsmidx}
\usepackage{stmaryrd}
%\usepackage{appendix}


\usepackage{apptools}
\usepackage{chngcntr}
%\AtAppendix{\counterwithin{thm}{chapter}}

%My Formating------------------------------------------------------

%\begin{comment}
\addtolength{\oddsidemargin}{-1.5cm}
\addtolength{\evensidemargin}{-1.5cm}
\addtolength{\textwidth}{3cm}
\addtolength{\topmargin}{-1cm}
\addtolength{\textheight}{1.3cm}
%\end{comment}

%Theorems and etc --------------------------------

\newtheorem{thm}{Theorem}[section]
\newtheorem{prop}[thm]{Proposition}
\newtheorem{lem}[thm]{Lemma}
\newtheorem{cor}[thm]{Corollary}
\newtheorem{conj}[thm]{Conjecture}
\newtheorem{claim}{Claim}

\theoremstyle{definition}
\newtheorem{defn}[thm]{Definition}
\newtheorem{notn}[thm]{Notation}
\newtheorem{example}[thm]{Example}
\newtheorem{exercise}{Exercise}[section]
\newtheorem{assumptions}[thm]{Assumptions}

\theoremstyle{remark}
\newtheorem{rem}[thm]{Remark}
\newtheorem{rems}[thm]{Remarks}


% -------------------------------------------------------------------

%\numberwithin{section}{chapter}
\numberwithin{equation}{section}
\numberwithin{thm}{section}


\makeatletter 
\newcommand\mynobreakpar{\vspace{0.02in}\par\nobreak\@afterheading}  
\makeatother

\interfootnotelinepenalty=10000 % Not to break up footnotes

\makeatletter
\@addtoreset{footnote}{section}
\makeatother

\usepackage{xcolor}
\hypersetup{
	colorlinks,
	linkcolor={red!50!black},
	citecolor={blue!50!black},
	urlcolor={blue!80!black}
}

\AtBeginDocument{%
	\def\MR#1{}
}


% Math Operators----------------------

\DeclareMathOperator{\sgn}{sgn}
\DeclareMathOperator{\dist}{dist}
\DeclareMathOperator{\Nul}{Nul}
\DeclareMathOperator{\Hom}{Hom}
\DeclareMathOperator{\End}{End}
\DeclareMathOperator{\Aut}{Aut}
\DeclareMathOperator{\rank}{rank}
\DeclareMathOperator{\ind}{ind}
\DeclareMathOperator{\Ht}{ht}
\DeclareMathOperator{\Frob}{Frob}
\DeclareMathOperator{\Fr}{Fr}
\DeclareMathOperator{\Spec}{Spec}
\DeclareMathOperator{\spec}{spec}
\DeclareMathOperator{\inv}{inv}
\DeclareMathOperator{\Lie}{Lie}
\DeclareMathOperator{\chr}{char}
\DeclareMathOperator{\re}{Re}
\DeclareMathOperator{\im}{Im}
\DeclareMathOperator{\Mat}{Mat}
\DeclareMathOperator{\ord}{ord}
\DeclareMathOperator{\Res}{Res}
\DeclareMathOperator{\Vol}{Vol}
\DeclareMathOperator{\Gal}{Gal}
\DeclareMathOperator{\GL}{GL}
\DeclareMathOperator{\PGL}{PGL}
\DeclareMathOperator{\SL}{SL}
\DeclareMathOperator{\coker}{coker}
\DeclareMathOperator{\Fix}{Fix}
\DeclareMathOperator{\Div}{Div}
\DeclareMathOperator{\Pic}{Pic}
\DeclareMathOperator{\Cl}{Cl}
\DeclareMathOperator{\Alb}{Alb}
\DeclareMathOperator{\Tr}{Tr}
\DeclareMathOperator{\Stab}{Stab}
\DeclareMathOperator{\diag}{diag}
\DeclareMathOperator{\Span}{span} 
\DeclareMathOperator{\Nr}{Nr}
\DeclareMathOperator{\rnr}{nr}
\DeclareMathOperator{\rtr}{tr}
\DeclareMathOperator{\disc}{disc}
\DeclareMathOperator{\Ind}{Ind}
\DeclareMathOperator{\Ann}{Ann}
\DeclareMathOperator{\NP}{NP}
\DeclareMathOperator{\Mtr}{Mat}
\DeclareMathOperator{\CM}{CM}
\DeclareMathOperator{\Reg}{Reg}
\DeclareMathOperator{\Spl}{Spl}
\DeclareMathOperator{\Cent}{Cent}
\DeclareMathOperator{\Li}{Li}
\DeclareMathOperator{\Ram}{Ram}
\DeclareMathOperator{\lcm}{lcm}
\DeclareMathOperator{\gcrd}{gcrd}


\DeclareRobustCommand\longtwoheadrightarrow
{\relbar\joinrel\twoheadrightarrow}

\newcommand{\free}{\mathrm{free}}
\newcommand{\sep}{\mathrm{sep}}
\newcommand{\alg}{\mathrm{alg}}
\newcommand{\id}{\mathrm{id}}
\newcommand{\an}{\mathrm{an}}
\newcommand{\wild}{\mathrm{w}}
\newcommand{\tame}{\mathrm{t}}
\newcommand{\ab}{\mathrm{ab}}
\newcommand{\nr}{\mathrm{nr}}
\newcommand{\rig}{\mathrm{rig}}
\newcommand{\tor}{\mathrm{tor}}
\newcommand{\new}{\mathrm{new}}
\newcommand{\old}{\mathrm{old}}
\newcommand{\df}{\mathrm{def}}
\newcommand{\red}{\mathrm{red}}
\newcommand{\modr}{\ \mathrm{mod}_r\ }

\newcommand{\me}{\mathrm{e}} % if one wants to get fancy with exp

\newcommand{\et}{\mathrm{et}}
\newcommand{\un}{\mathrm{nr}}
\newcommand{\ev}{\mathrm{ev}}
\newcommand{\rH}{\mathrm{H}}

% Frankfurt ---------------------------

\newcommand{\fA}{\mathfrak{A}}
\newcommand{\fB}{\mathfrak{B}}
\newcommand{\fC}{\mathfrak{C}}
\newcommand{\fD}{\mathfrak{D}}
\newcommand{\fE}{\mathfrak{E}}
\newcommand{\fF}{\mathfrak{F}}
\newcommand{\fG}{\mathfrak{G}}
\newcommand{\fH}{\mathfrak{H}}
\newcommand{\fI}{\mathfrak{I}}
\newcommand{\fJ}{\mathfrak{J}}
\newcommand{\fK}{\mathfrak{K}}
\newcommand{\fL}{\mathfrak{L}}
\newcommand{\fM}{\mathfrak{M}}
\newcommand{\fN}{\mathfrak{N}}
\newcommand{\fO}{\mathfrak{O}}
\newcommand{\fP}{\mathfrak{P}}
\newcommand{\fQ}{\mathfrak{Q}}
\newcommand{\fR}{\mathfrak{R}}
\newcommand{\fS}{\mathfrak{S}}
\newcommand{\fT}{\mathfrak{T}}
\newcommand{\fU}{\mathfrak{U}}
\newcommand{\fV}{\mathfrak{V}}
\newcommand{\fW}{\mathfrak{W}}
\newcommand{\fX}{\mathfrak{X}}
\newcommand{\fY}{\mathfrak{Y}}
\newcommand{\fZ}{\mathfrak{Z}}

\newcommand{\fa}{\mathfrak{a}}
\newcommand{\fb}{\mathfrak{b}}
\newcommand{\fc}{\mathfrak{c}}
\newcommand{\fd}{\mathfrak{d}}
\newcommand{\fe}{\mathfrak{e}}
\newcommand{\ff}{\mathfrak{f}}
\newcommand{\fg}{\mathfrak{g}}
\newcommand{\fh}{\mathfrak{h}}
\newcommand{\fri}{\mathfrak{i}}
\newcommand{\fj}{\mathfrak{j}}
\newcommand{\fk}{\mathfrak{k}}
\newcommand{\fl}{\mathfrak{l}}
\newcommand{\fm}{\mathfrak{m}}
\newcommand{\fn}{\mathfrak{n}}
\newcommand{\fo}{\mathfrak{o}}
\newcommand{\fp}{\mathfrak{p}}
\newcommand{\fq}{\mathfrak{q}}
\newcommand{\fr}{\mathfrak{r}}
\newcommand{\fs}{\mathfrak{s}}
\newcommand{\ft}{\mathfrak{t}}
\newcommand{\fu}{\mathfrak{u}}
\newcommand{\fv}{\mathfrak{v}}
\newcommand{\fw}{\mathfrak{w}}
\newcommand{\fx}{\mathfrak{x}}
\newcommand{\fy}{\mathfrak{y}}
\newcommand{\fz}{\mathfrak{z}}

% Calligraphic--------------------------

\newcommand{\cA}{\mathcal{A}}
\newcommand{\cB}{\mathcal{B}}
\newcommand{\cC}{\mathcal{C}}
\renewcommand{\cD}{\mathcal{D}}
\newcommand{\cE}{\mathcal{E}}
\newcommand{\cF}{\mathcal{F}}
\newcommand{\cG}{\mathcal{G}}
\renewcommand{\cH}{\mathcal{H}}
\newcommand{\cI}{\mathcal{I}}
\newcommand{\cJ}{\mathcal{J}}
\newcommand{\cK}{\mathcal{K}}
\renewcommand{\cL}{\mathcal{L}}
\newcommand{\cM}{\mathcal{M}}
\newcommand{\cN}{\mathcal{N}}
\newcommand{\cO}{\mathcal{O}}
\newcommand{\cP}{\mathcal{P}}
\newcommand{\cQ}{\mathcal{Q}}
\renewcommand{\cR}{\mathcal{R}}
\newcommand{\cS}{\mathcal{S}}
\newcommand{\cT}{\mathcal{T}}
\newcommand{\cU}{\mathcal{U}}
\newcommand{\cV}{\mathcal{V}}
\newcommand{\cW}{\mathcal{W}}
\newcommand{\cX}{\mathcal{X}}
\newcommand{\cY}{\mathcal{Y}}
\newcommand{\cZ}{\mathcal{Z}}

%Script-------------------------

\newcommand{\sA}{\mathscr{A}}
\newcommand{\sB}{\mathscr{B}}
\newcommand{\sC}{\mathscr{C}}
\newcommand{\sD}{\mathscr{D}}
\newcommand{\sE}{\mathscr{E}}
\newcommand{\sF}{\mathscr{F}}
\newcommand{\sG}{\mathscr{G}}
\newcommand{\sH}{\mathscr{H}}
\newcommand{\sI}{\mathscr{I}}
\newcommand{\sJ}{\mathscr{J}}
\newcommand{\sK}{\mathscr{K}}
\newcommand{\sL}{\mathscr{L}}
\newcommand{\sM}{\mathscr{M}}
\newcommand{\sN}{\mathscr{N}}
\newcommand{\sO}{\mathscr{O}}
\newcommand{\sP}{\mathscr{P}}
\newcommand{\sQ}{\mathscr{Q}}
\newcommand{\sR}{\mathscr{R}}
\newcommand{\sS}{\mathscr{S}}
\newcommand{\sT}{\mathscr{T}}
\newcommand{\sU}{\mathscr{U}}
\newcommand{\sV}{\mathscr{V}}
\newcommand{\sW}{\mathscr{W}}
\newcommand{\sX}{\mathscr{X}}
\newcommand{\sY}{\mathscr{Y}}
\newcommand{\sZ}{\mathscr{Z}}

%BB------------------------------

\newcommand{\A}{\mathbb{A}}
\newcommand{\B}{\mathbb{B}}
\newcommand{\C}{\mathbb{C}}
\newcommand{\D}{\mathbb{D}}
\newcommand{\E}{\mathbb{E}}
\newcommand{\F}{\mathbb{F}}
\newcommand{\gm}{\mathbb{G}}
\renewcommand{\H}{\mathbb{H}}
\newcommand{\I}{\mathbb{I}}
\newcommand{\J}{\mathbb{J}}
\newcommand{\K}{\mathbb{K}}
\renewcommand{\L}{\mathbb{L}}
\newcommand{\M}{\mathbb{M}}
\newcommand{\N}{\mathbb{N}}
\renewcommand{\O}{\mathbb{O}}
\newcommand{\p}{\mathbb{P}}
\newcommand{\Q}{\mathbb{Q}}
\newcommand{\R}{\mathbb{R}}
\newcommand{\s}{\mathbb{S}}
\newcommand{\T}{\mathbb{T}}
\newcommand{\U}{\mathbb{U}}
\newcommand{\V}{\mathbb{V}}
\newcommand{\W}{\mathbb{W}}
\newcommand{\X}{\mathbb{X}}
\newcommand{\Y}{\mathbb{Y}}
\newcommand{\Z}{\mathbb{Z}}


%Bold------------------------------------

\newcommand{\bone}{{\bm{1}}}
\newcommand{\beps}{{\bm{\epsilon}}}
\newcommand{\bom}{\boldsymbol{m}}
\newcommand{\boz}{\boldsymbol{z}}
\newcommand{\bou}{\boldsymbol{u}}
\newcommand{\bow}{\boldsymbol{w}}
\newcommand{\boe}{\boldsymbol{e}}
\newcommand{\bov}{\boldsymbol{v}}
\newcommand{\bok}{\boldsymbol{k}}
\newcommand{\boi}{\mathbf{i}}
\newcommand{\bof}{\mathbf{f}}



% Miscellaneous Symbols---------------------------

\newcommand{\tF}{\widetilde{F}}
\newcommand{\hF}{\widehat{F}}
\newcommand{\hK}{\widehat{K}}
\newcommand{\hR}{\widehat{R}}
\newcommand{\hM}{\widehat{\mathcal{M}}}
\newcommand{\tinf}{\widetilde{\infty}}
\newcommand{\tfp}{\widetilde{\fp}}
\newcommand{\hphi}{\widehat{\phi}}

\newcommand{\eps}{\varepsilon}
\newcommand{\G}{\Gamma}
\newcommand{\oG}{\overline{\Gamma}}
\newcommand{\To}{\longrightarrow}
\newcommand{\bs}{\setminus}
\newcommand{\Fi}{F_\infty}
\newcommand{\Ci}{\C_\infty}
\newcommand{\Ki}{K_\infty}
\newcommand{\oF}{\overline{\F}}
\newcommand{\oK}{\overline{K}}
\newcommand{\oM}{\overline{M}}

\newcommand{\La}{\Lambda}
\newcommand{\la}{\lambda}

\newcommand{\norm}[1]{\left\Vert#1\right\Vert}

\newcommand{\twist}[1]{#1\!\left\{\tau\right\}}
\newcommand{\btwist}[1]{#1\!\left\{\tau^{-1}\right\}}
\newcommand{\twistt}[2]{#1\!\left\{#2\right\}}
\newcommand{\atwist}[2]{#1\langle#2\rangle}

\newcommand{\twistdvr}[1]{#1\!\left\{\mkern-4mu\left\{\tau\right\}\mkern-4mu\right\}}
\newcommand{\atwistdvr}[2]{#1\!\left\langle\mkern-4mu\left\langle#2\right\rangle\mkern-4mu\right\rangle}
\newcommand{\ls}[2]{#1\!\left(\mkern-4mu\left(#2\right)\mkern-4mu\right)}
\newcommand{\dvr}[2]{#1\!\left\llbracket #2\right\rrbracket}

\newcommand{\Leg}[2]{\genfrac{(}{)}{}{}{#1}{#2}}
\newcommand{\ReS}[3]{\genfrac{(}{)}{}{}{#1}{#2}_{\mkern-5mu #3}}
\newcommand{\Mod}[1]{\ (\mathrm{mod}\ #1)}
\newcommand{\abs}[1]{\left|#1\right|}

\newcommand{\Cbin}[2]{\left\{ {#1\atop #2} \right\}}


\newcommand{\prodprime}{\sideset{}{'}\prod}
\newcommand{\sumprime}{\sideset{}{'}\sum}
\newcommand{\lr}[1]{\langle#1\rangle}

\newcommand{\circo}{~\raisebox{1pt}{\tikz \draw[line width=0.6pt] circle(1.1pt);}~}


\title{Rank-metric codes from Drinfeld modules}

\author{Giacomo Micheli}
\address{Department of Mathematics \& Statistics, The University of South Florida, Florida, United States of America}
\email{gmicheli@usf.edu}

\author{Mihran Papikian}
\address{Department of Mathematics, Pennsylvania State University, University Park, Pennsylvania, United States of America}
\email{papikian@psu.edu}

\thanks{The first author was supported by NSF CAREER grant 2338424. The second author was supported in part by the Simons Foundation, award number MPS-TSM-00008093.} 


\subjclass[2020]{11G09, 12K10, 11T71, 16S36}
\keywords{Rank metric codes, Drinfeld modules, twisted polynomial ring}





% ----------------------------------------------------------------------
\begin{document}
	


\begin{abstract}
We establish a connection between Drinfeld modules and rank-metric codes, focusing  
on the case of semifield codes. 
Our method constructs rank-metric codes from linear subspaces of endomorphisms of a 
Drinfeld module acting on torsion submodules. We show that Sheekey's construction \cite{Sheekey} 
fits naturally into this framework, yielding a short conceptual proof of one of his main results. 
We then give a new construction of infinite families of semifield codes arising from 
Drinfeld modules defined over finite fields.
\end{abstract}

%\begin{color}{red}
%Since we cannot prove novelty, I would push on the fact that our polynomials have high degree %or something similar. What makes them special with respect to Sheekey?
%\end{color}

\maketitle
\section{Introduction} 

\subsection{Rank-metric codes and semifields}\label{ss1.1} Error-correcting codes are used when communication is over a channel in which 
errors may occur. This requires a set equipped with a distance function and a subset of allowed codewords; if the number of errors 
is assumed to be small, then a received message is decoded to be the nearest valid codeword. 

The most well-known and widely-used examples are the codes with the Hamming metric: the codewords are taken from an $n$-dimensional vector 
space $\F_q^n$ over a finite field $\F_q$ with $q$ elements (usually $q=2$), and the distance between two vectors in $\F_q^n$ 
is defined as the number of positions in which they differ.  

In rank-metric coding, the codewords are instead taken from the set $\Mat_{m, n}(\F_q)$  of $m\times n$-matrices ($m\leq n$), with 
the distance between two matrices defined as the rank of their difference 
$$
\dist(X, Y) = \rank (X-Y). 
$$
Rank-metric codes have seen renewed interest in recent years due to their applications in network coding, distributed
data storage,  
and cryptography; cf. \cite{SKK}, \cite{bartz2022}, \cite{SheekeySurvey}. 

A rank-metric code is a subset $\cC \subset \Mat_{m, n}(\F_q)$ with a minimum distance defined by 
$$
\dist(\cC)=\min \{\dist(X, Y)\mid X, Y\in \cC, X\neq Y\}. 
$$
Such $\cC$ satisfies the  
 \textit{Singleton-type bound}: if $d=\dist(\cC)$, then 
\begin{equation}\label{eqSingleton}
	\abs{\cC} \leq q^{n(m-d+1)}. 
\end{equation}
A code obtaining the bound \eqref{eqSingleton} is said to be a \textit{Maximum Rank Distance code}, or \textit{MRD code} for short. It is know that 
MRD codes exist for every finite field $\F_q$ and any choice of parameters $d\leq m\leq n$; see  \cite{Delsarte}. An MRD code 
can correct up to $\lfloor (d-1)/2\rfloor$ errors. 



We will assume for the rest of this article that MRD codes that we consider are $\F_q$-linear, i.e., $\cC$ 
is an $\F_q$-vector subspace of $\Mat_{m, n}(\F_q)$. 
In the special case when $d=m=n$, an $\F_q$-linear MRD code is equivalent to a 
\textit{semifield}---a finite dimensional $\F_q$-algebra with unity, in which multiplication is not 
necessarily commutative or associative; see \cite[Prop. 1]{Sheekey}.  
To distinguish such MRD codes, 
we will call them \textit{semifield codes}, although this is not a standard terminology. 

To see the connection between semifields and MRD codes, let $\cS$ be a semifield of dimension $n$ over $\F_q$.  
Then the right multiplication by a non-zero element $\alpha$ on $\cS$ defines an 
invertible linear operator $T_\alpha\colon \cS \to \cS$, $\beta\mapsto \beta\alpha$. After fixing a basis of $\cS$ over $\F_q$, we get 
an injection $\cS\to \Mat_n(\F_q)$, whose image is a semifield code. 
For example, this can be applied to $\cS=k$, where 
$$k\colonequals \F_{q^n}$$ is the finite field with $q^n$ elements, considered as an $\F_q$-vector space.  

\begin{example}
A more interesting example of a finite semifield arises from the \textit{twisted polynomial ring} 
$$
R\colonequals \twist{k}=\left\{\sum_{i=0}^n a_i \tau^i \mid n\geq 0, a_i\in k\right\}, 
$$ 
where the addition is the usual addition of polynomials, but 
the multiplication is subject to the following commutation rule $$\tau a=a^q\tau \quad \text{for all} \quad a\in k.$$ It is clear that $R$ is a (noncommutative) ring without zero-divisors. It is also easy to show that $R$ possesses a right (and left) division algorithm. % (see \cite[Thm. 3.1.13]{PapikianGTM}). 
Let $P\in R$ be a monic  irreducible element of degree $s$. For any $f\in R$, there are unique $g, e\in R$ such that 
$$
f=gP+e \quad \text{with } \deg(e)<s \text{ or }e=0. 
$$
We denote the residue $e=f\modr P$. With this notation, the set 
$$
\cS=\{f\in R\mid \deg(f)\leq s-1\},  
$$
with the addition from $R$ and multiplication 
$a\circ b = ab\modr P$ 
is a semifield of order $q^{ns}$. The MRD codes that one obtains from such semifields are called \textit{Gabidulin codes}; cf. \cite{Sheekey}.   
\end{example}

Elliptic curves have enjoyed great success in cryptographic applications. 
While Drinfeld modules---despite their many formal similarities to elliptic curves---are not suitable for cryptographic applications due to their semilinear nature \cite{Scanlon}, this same semilinear structure turns out to be ideally suited for constructing rank-metric codes. 
In this paper, we explain how the theory of Drinfeld modules can be used to construct rank-metric codes. 
The Drinfeld module perspective may situate various constructions of such codes, 
e.g. \cite{Sheekey}, \cite{TZ},  \cite{LSS}, \cite{GTLNS}, within a unified algebraic framework, 
opening the door to powerful tools from the theory of Drinfeld modules into coding theory.

%------------------------------------------------------

\subsection{Drinfeld modules}\label{sIntroDM} 
We denote $A=\F_q[T]$, the polynomial ring in indeterminate $T$ with coefficients in $\F_q$. 
Given a nonzero ideal $\fn\lhd A$, with slight abuse of notation, we denote by the same symbol  
the unique monic generator of $\fn$. The \textit{primes} of $A$ are the maximal ideals of $A$. 
Given a prime $\fp\lhd A$, we denote $\F_\fp=A/\fp$. 

A \textit{Drinfeld module} of rank $r\geq 1$ over $k$ is an $\F_q$-algebra homomorphism 
\[\phi\colon A\to R,\quad  a\mapsto \phi_a, \] 
such that $\phi_T=g_0+g_1\tau+\cdots+g_r\tau^r$ for some $g_0, \dots, g_r\in k$ with $g_r\neq 0$. The \textit{$A$-characteristic of $\phi$}, 
denoted $\chr_A(\phi)$, is the kernel of the homormorphism $A\to k$ defined by $T\mapsto g_0$. The Drinfeld module $\phi$ makes the algebraic 
closure $\bar{k}$ into an $A$-module, denoted ${^\phi}\overline{k}$; see Section \ref{sDM} for the details. The elements of ${^\phi}\overline{k}$ 
annihilated by $a\in A$ form a finite $A$-module denoted $\phi[a]$; if $\chr_A(\phi)\nmid a$, then $\phi[a]\cong (A/aA)^r$. 

Define 
the \textit{endomorphism ring of $\phi$} as 
$$
\End(\phi)=\left\{u\in R\mid u\phi_T=\phi_Tu\right\}. 
$$
$\End(\phi)$ is an $A$-algebra, which is known to be a free $A$-module of rank $\leq r^2$. Every $u\in \End(\phi)$ acts $A$-linearly on $\phi[a]$. 
Choosing a prime $\fp\neq \chr_A(\phi)$, we obtain a homomorphism 
$$
\iota_\fp\colon \End(\phi)\To \End_A(\phi[\fp])\cong \Mat_r(\F_\fp). 
$$

\vspace{0.1in}
\noindent 
\textbf{\underline{The recipe}.} Our construction of semifield codes relies on three main ingredients:
\begin{itemize}
\item A Drinfeld module $\phi$ of rank $r$.
\item An $\F_q$-linear subspace $\cM \subset \End(\phi)$.
\item A prime $\fp \in A$, distinct from the $A$-characteristic of $\phi$, such that every nonzero $u \in \cM$ acts \textit{invertibly} on $\phi[\fp]$, i.e., as an invertible linear transformation. 
\end{itemize}
Given this, we obtain an embedding
\[
\iota_\fp \colon \cM \longhookrightarrow \End_{\F_\fp}(\phi[\fp]) \cong \Mat_r(\F_\fp), 
\]
where each nonzero $\iota_\fp(u)$ has full rank. If additionally 
\[\dim_{\F_q} \cM = r \cdot \deg(\fp),\] then the image $\iota_\fp(\cM)$ is a semifield code.

\vspace{0.1in}

The flexibility of choosing $\phi$ and $\cM$ makes this construction very general. 
The advantage of this perspective is that it provides access to the rich theory of 
Drinfeld modules. The central challenge is to find a suitable $\fp$ satisfying the invertibility condition; here one can leverage various tools from the theory of Drinfeld modules, such as the equality between the $\tau$-degree of an endomorphism and the $T$-degree of its determinant on the associated Tate module. We focus on the semifield case, though this strategy extends to MRD codes by relaxing the invertibility requirement.

\begin{example}\label{BasicExample}
We demonstrate the idea with a simple example. 
Define $\phi$ by $\phi_T=\tau^n$. The rank of $\phi$ is $n$, $\chr_A(\phi)=T$, and $\End(\phi)=R$. 
Let $\fp\neq T$ be a prime of degree $s$ and let 
$$
\cM=\{u\in R\mid \deg(u)\leq s-1\}. 
$$
We claim that every $0\neq u\in \cM$ acts invertibly on $\phi[\fp]$. We argue by contradiction. Assume $0\neq \beta\in \phi[\fp]$ 
is annihilated by $u$. Then the $A$-submodule of $\phi[\fp]$ generated by $\beta$ has size $q^s$, yet $\deg(u)<s$ implies $\#\ker(u)\leq q^{s-1}$, a contradiction. 
Since $\dim_{\F_q}\cM=ns=\deg(\fp)\cdot n$, the image $\iota_\fp(\cM)$ is a semifield code.
\end{example}

\subsection{Main results} 
%The first construction is in fact due to Sheekey \cite{Sheekey}, who works in \textit{loc. cit.} in the setting of linearized polynomials, without using 
%the theory of Drinfeld modules. 
In Section \ref{secSheekey}, we give a reinterpretation of Sheekey's construction 
\cite{Sheekey} in terms of Drinfeld modules and a short proof 
of one of the main results in \cite{Sheekey}. 
Sheekey's construction has been generalized in \cite{TZ} and \cite{LSS} using the 
methods in \cite{Sheekey}, which rely on direct computations in skew-polynomial rings 
and their quotients. 

In our Drinfeld module reinterpretation of Sheekey's construction, as in Example \ref{BasicExample}, one takes $\phi_T=\tau^n$, but now 
$$
\cM \colonequals \{u_0+u_1\tau\cdots+u_s\tau^s\in R\mid  u_s=\eta u_0\}
$$
for a fixed $\eta\in k$. Let $\fp\neq T$ be a monic prime in $A$ of degree $s$ with constant coefficient $\fp_0$. We show that if 
$\Nr_{k/\F_q}(\eta)\cdot (-1)^{s(n-1)}\fp_0\neq 1$, then the image of $\cM$ in $\End_{\F_\fp}(\phi[\fp])$ is a semifield code. The proof 
uses the duality between $u \in \cM$ 
being an endomorphism of $\phi$ and $\phi_T$ being an endomorphism of the auxiliary 
Drinfeld module $\psi$ defined by $\psi_T=u$: this converts invertibility questions into questions about 
characteristic polynomials of endomorphisms, for which standard tools from the theory of Drinfeld 
modules are available. Note that choosing $\eta=0$ gives Example \ref{BasicExample}. The interesting fact about Sheekey's construction 
is that the semifield codes thus constructed in some cases correspond to genuinely new families of semifields, with 
parameters for which no examples were known before. 


In Section \ref{s2C}, we extend the setup as follows: 
Let $\ell$ and $s$ be integers such that $n=\lcm(\ell, s)$. We choose $\phi$ so that $\phi_T\in \twistt{\F_{q^s}}{\tau^\ell}\subset \twist{k}$ 
and $\deg(\phi_T)=r\ell$. In this case, $\twistt{\F_{q^\ell}}{\tau^s}\subset \End(\phi)$. Let $\fp$ be a prime different from 
$\chr_A(\phi)$ and denote $d=\deg_T(\fp)$. We take 
$$
\cM=\left\{\sum_{i=0}^{rd-1} b_i\tau^{s i}\mid b_i\in \F_{q^\ell}\right\}. 
$$

Unlike the constructions in \cite{Sheekey, TZ, LSS, GTLNS}, which all take $\phi_T = \tau^n$, 
our construction uses Drinfeld modules with smaller endomorphism rings, exploiting 
the structure of division subalgebras of $k(\tau)$.
Under some technical assumptions on the reduction modulo $\fp$ of the minimal polynomial $m_\phi(x)$ of the Frobenius endomorphism $\pi=\tau^n\in \End(\phi)$, 
we show that the image of $\cM$ in $\End_{\F_\fp}(\phi[\fp])$ is a semifield code. Here again the duality between $u\in \cM$ being an 
endomorphism of $\phi$ and $\phi_T$ being an endomorphism of $\psi$ defined by $\psi_T=u$ comes into play. Next, 
using the Chebotarev density theorem, we show that one can always choose $\phi$ and $\fp$ so that the assumptions on $m_\phi(x)$ 
are satisfied. Finally, we compute the nuclear parameters of our code and illustrate the construction with an 
explicit example. While the nuclear parameters we obtain do not immediately 
distinguish our codes from known families, the construction is conceptually new  
and may yield new semifields for appropriate choices of parameters.

% ----------------------------------------------------------------------

\section{Preliminaries on Drinfeld modules}\label{sDM}  In this section we recall some basic facts 
from the theory of Drinfeld modules that will be used later in the paper. Most of the proofs can be found in \cite{PapikianGTM}. 

\subsection{Definition} Let $F=\F_q(T)$ be the fraction field of $A=\F_q[T]$. For $0\neq a\in A$, let $\deg(a)\in \Z_{\geq 0}$ be the usual degree 
of $a$ as a polynomial in $T$. 
%(When degrees of two or more polynomial rings appear in the same formula we sometimes denote $\deg(a)$ by $\deg_T(a)$ for clarity.) 
%The degree map naturally extends to $F$ by $\deg(a/b)=\deg(a)-\deg(b)$. The map $\ord_\infty= -\deg\colon F\to \Z\cup \{-\infty\}$ is a valuation and the corresponding place of $F$ is denoted by $\infty$. We normalize the absolute value $\abs{\cdot}$  on $F$ corresponding to $\ord_\infty$by $\abs{a}=q^{\ord_\infty(a)}$. Note that for $a\in A$ we have $\abs{a}=\# (A/aA)$. 

Let $K$ be a field containing $\F_q$ as a subfield. Let 
$\twist{K}$ be the twisted polynomial ring already discussed in $\S$\ref{ss1.1} in the case when $K=\F_{q^n}$ (so the commutation rule is $\tau a=a^q\tau$ for all $a\in K$). 
There is a homomorphism \[ \partial\colon \twist{K}\to K, \quad \sum_{i=0}^n a_i \tau^i \mapsto a_0,\] called the \textit{derivative}. 
For $f=a_h\tau^h+\cdots +a_n\tau^n$ with $0\leq h\leq n$, $a_h\neq 0$ and $a_n\neq 0$, we define the \textit{height} of $f$ as $\Ht(f)=h$, and the \textit{degree} of $f$ as $\deg_\tau(f)=n$. 
The map $$f=\sum_{i=0}^n a_i \tau^i \longmapsto f(x)=\sum_{i=0}^n a_i x^{q^i}$$ gives a ring isomorphism between $\twist{K}$ 
and the ring $\atwist{K}{x}$ of $\F_q$-linear polynomial, where on $\atwist{K}{x}$ the multiplication is defined via the substitution $f_1\circ f_2=f_1(f_2(x))$. 
%Note that on $\atwist{K}{x}$ the derivative $\partial$ is the usual formal derivative of polynomials with respect to the variable $x$, i.e., $\partial(f)=\frac{d}{dx}f(x)$. 

%We say that $K$ is an \textit{$A$-field} if $K$ is equipped with with an $\F_q$-algebra homomorphism $\gamma\colon A\to K$. The $A$-characteristic of $K$, 
%denoted $\chr_A(K)$, is the kernel of $\gamma$; thus, either $\chr_A(K)=0$ or $\chr_A(K)=\fp$ is a prime of $A$. 

A \textit{Drinfeld module} of rank $r\geq 1$ over $K$ is an $\F_q$-algebra homomorphism 
\begin{align*}
	\phi\colon A &\To \twist{K}, \\ 
	a &\longmapsto \phi_a=g_0(a)+g_1(a)\tau+\cdots +g_n(a)\tau^n,
\end{align*}
such that for $a\neq 0$ we have $n=\deg(a) r$ and $g_n(a)\neq 0$. Note that $\phi$ is always injective, so gives an embedding of $A$ 
into the non-commutative ring $\twist{K}$. Moreover, $\phi$ is uniquely determined by $\phi_T$, so to define a Drinfeld module over $K$ 
one simply needs to choose $g_0, g_1, \dots, g_r\in K$ such that $g_r\neq 0$ and put $\phi_T=g_0+g_1\tau+\cdots+g_r\tau^r$. 
We define a homomorphism $\gamma\colon A\to K$ via $a\mapsto \partial \phi_a$. The \textit{$A$-characteristic of $\phi$} is 
$\chr_A(\phi)\colonequals \ker(\gamma)$; note that $\chr_A(\phi)$ is either $0$ or a prime of $A$. 

\subsection{Torsion module and endomorphisms}\label{secTM} 
Let $\phi\colon A\to \twist{K}$ be a Drinfeld module over $K$ of rank $r$. Through $\phi$, the algebraic closure $\oK$ of $K$ acquires an $A$-module structure, where $a\in A$ acts on $\beta\in \oK$ by $a\ast \beta=\phi_a(\beta)$, i.e., we substitute $\beta$ into the 
polynomial $\phi_a(x)$. Denote this module by ${^\phi}{\oK}$. One is interested in its torsion submodule. More precisely, given $0\neq a\in A$, 
the \textit{$a$-torsion submodule} of ${^\phi}{\oK}$, denoted $\phi[a]$,  is the set of roots of the polynomial $\phi_a(x)$. It is easy to check that $\phi[a]$ 
is indeed an $A$-submodule of ${^\phi}{\oK}$, i.e., is invariant under the action of any $b\in A$, where $b$ acts via $\phi_b$. 
If $\chr_A(\phi)$ does not divide $a$, then $\phi[a]\cong (A/aA)^r$; see \cite[Cor. 3.5.3]{PapikianGTM}. 

Let $\fp\lhd A$ be a prime different from $\chr_A(\phi)$. The finite $A$-modules $\phi[\fp^n]\cong (A/\fp^n)^r$, $n\geq 1$, form a projective system 
with transition maps $\phi[\fp^{n+1}]\to \phi[\fp^n]$, $\alpha\mapsto \phi_\fp(\alpha)$. Taking the inverse limit, one obtains the \textit{$\fp$-adic Tate module} 
$T_\fp(\phi)\cong A_\fp^r$ of $\phi$, where $A_\fp$ is the ring of integers of the completion of $F$ with respect to the absolute value arising from the $\fp$-adic valuation on $F$. 

The \textit{endomorphisms of $\phi$} defined over $K$ are the elements of $\twist{K}$ which induce endomorphisms of ${^\phi}\oK$. More explicitly,
\begin{align*}
\End(\phi) &=\{u\in \twist{K}\mid u\phi_a=\phi_au\text{ for all }a\in A\} \\ 
& = \{u\in \twist{K}\mid u\phi_T=\phi_Tu\}. 
\end{align*}
Obviously, the image $\phi(A)$ of $A$ under $\phi$ is in $\End(\phi)$. In fact, it is easy to check that $\End(\phi)$ is an $A$-algebra, where we 
identify $\phi(A)$ with $A$. Moreover, one can show 
that $\End(\phi)$ is a free $A$-module of rank $\leq r^2$; see \cite[Thm. 3.4.1]{PapikianGTM}. 

Any $u\in \End(\phi)$ induces an $A$-linear transformation on $\phi[a]$, and also an $A_\fp$-linear transformation on $T_\fp(\phi)$. 
Let $P_{\phi, u}(x)\in A_\fp[x]$ denote the characteristic polynomial of this latter transformation. It turns out that the coefficients of $P_{\phi, u}(x)$ are 
actually in $A$ and do not depend on $\fp$; see \cite[Thm. 3.6.6]{PapikianGTM}. The characteristic polynomial of $u$ acting on 
on the $\F_\fp$-vector space $\phi[\fp]\cong \F_\fp^r$ is the reduction modulo $\fp$ of $P_{\phi, u}(x)$. In particular, $u$ acts \textit{invertibly} on $\phi[\fp]$ 
 (i.e., as an invertible linear transformation) if and only if $P_{\phi, u}(0)$ is not divisible by $\fp$.  
 
 One of the key technical questions arising in this 
 paper is to show that certain endomorphisms act invertibly on $\phi[\fp]$ for an appropriately chosen $\fp$. More precisely, 
 we choose a specific $\F_q$-vector subspace $\cM\subset \End(\phi)$--\textit{the message space}--and show that every $0\neq u\in \cM$
acts invertibly on $\phi[\fp]$. In that situation, $\cM$ defines a semifield code in $\End_{\F_\fp}(\phi[\fp])\cong \Mat_r(\F_\fp)$ if 
\begin{equation}\label{eqSFdim}
\dim_{\F_q}\cM=r\cdot \deg(\fp). 
\end{equation}
 %-------------------------
 
 \subsection{Anderson motive of Drinfeld module}\label{secMotive} Let $\phi\colon A\to 
 \twist{K}$ be a Drinfeld module over $K$ of rank $r$ defined by 
 $$
 \phi_T=g_0+g_1\tau+\cdots+g_r\tau^r. 
 $$
 We associate with $\phi$ a left $K[T, \tau]$-module $M_\phi$, called the \textit{Anderson motive of $\phi$} \cite{Anderson}, by taking $M_\phi=\twist{K}$ and defining 
 \begin{align*}
 	u\circ m & = um \quad \text{for all $u\in \twist{K}$ and $m\in M_\phi$}, \\ 
 	a\circ m & = m\phi_a \quad \text{for all $a\in A$ and $m\in M_\phi$}, 
 \end{align*}
 where $um$ and $m\phi_a$ are multiplications in $\twist{K}$. By the proof of \cite[Lem. 3.4.4]{PapikianGTM}, $M_\phi$ 
 is freely generated over $K[T]$ by the elements  $\{1, \tau, \dots, \tau^{r-1}\}$; we call this set the \textit{standard basis} of $M_\phi$. 
 
 Let $u\in \End(\phi)$. Then $u$ induces an endomorphism of the $K[T, \tau]$-module $M_\phi$ by $m\mapsto mu$; see the proof of \cite[Prop. 3.4.5]{PapikianGTM}. 
  Let $U=(u_{i,j})_{0\leq i, j\leq r-1}\in \Mat_r(K[T])$ be the matrix 
 by which $u$ acts on $M_\phi$ as a free $K[T]$-module with respect to the standard basis, i.e., 
 $$
 \tau^i u = u_{0, i} + u_{1, i}\tau+\cdots +u_{r-1, i}\tau^{r-1}, \qquad 0\leq i\leq r-1. 
 $$
 By Theorem 3.6.6 and Proposition 3.6.7 in \cite{PapikianGTM}, 
 \begin{equation}\label{eqPu}
 P_{\phi, u}(x)=\det(xI_r-U),
 \end{equation}
 i.e., the characteristic polynomial of $u$ is equal to the characteristic polynomial of the matrix $U$. 
 In particular, $u$ acts invertibly on $\phi[\fp]$ if and only if $\fp$ does not divide $\det(U)$ in $A$. 
 
 For $f=f_0+f_1T+\cdots+f_nT^n\in K[T]$ and $i\geq 0$, denote $$f^{(q^i)} = f_0^{q^i}+f_1^{q^i}T+\cdots+f_n^{q^i}T^n.$$  
 For a matrix $S=(s_{ij})\in \Mat_{m, n}(K[T])$, put $S^{(q^i)}=\left(s_{ij}^{(q^i)}\right)$. 
 Let $$u=u_0+u_1\tau+\cdots+u_{r-1}\tau^{r-1}$$ be the expansion of $u$ in the standard basis. Then 
 \begin{align*}
 	\tau u & = u_0^{(q)}\tau+u_1^{(q)}\tau^2+\cdots+u_{r-1}^{(q)}\tau^r \\ 
 	& = u_0^{(q)}\tau+u_1^{(q)}\tau^2+\cdots+u_{r-1}^{(q)}(g_r^{-1}(T-g_0-g_1\tau-\cdots-g_{r-1}\tau^{r-1})) \\ 
 	& = u_{r-1}^{(q)}(g_r^{-1}(T-g_0)) + (u_0^{(q)}-u_{r-1}^{(q)}(g_r^{-1}g_1)) \tau+(u_1^{(q)}-u_{r-1}^{(q)}(g_r^{-1}g_2)) \tau^2+\cdots \\ &+
 	(u_{r-2}^{(q)}-u_{r-1}^{(q)}(g_r^{-1}g_{r-1})) \tau^{r-1}. 
 \end{align*}
 Put 
 \begin{equation}\label{eqSphi}
 	S_\phi=\begin{bmatrix} 
 		0 & 0 &\cdots & 0 &(T-g_0)/g_r \\ 
 		1 & 0 & \cdots & 0 &-g_1/g_r \\ 
 		\vdots & \vdots& \cdots &\vdots & \vdots \\ 
 		0 & 0 &\cdots & 1 & -g_{r-1}/g_r
 	\end{bmatrix}. 
 \end{equation}
 From the above computation one deduces that for any $1\leq i\leq r-1$ we have 
 \begin{align*}
 	\begin{bmatrix} u_{0,i}\\ \vdots \\ u_{r-1, i}\end{bmatrix} = 	S_\phi \begin{bmatrix} u_{0, i-1}\\ \vdots\\ u_{r-1, i-1}\end{bmatrix}^{(q)} & = 
 	S_\phi S_\phi^{(q)}  \begin{bmatrix} u_{0, i-2}\\ \vdots\\ u_{r-1, i-2}\end{bmatrix}^{(q^2)} = \cdots \\ &= 
 	S_\phi S_\phi^{(q)} \cdots  S_\phi^{(q^{i-1})} \begin{bmatrix} u_{0}\\ \vdots\\ u_{r-1}\end{bmatrix}^{(q^i)}. 
 \end{align*}

\begin{example}\label{ssExample} Here we construct a semifield code using the tools discussed earlier in this section. 

Let $r$ and $s$ be positive integers, let $n=\lcm(r, s)$, let $k=\F_{q^n}$, and let $t\in \F_{q^s}$. 
Let $\phi\colon A\to \twist{k}$ be the Drinfeld module defined by 
$$\phi_T=t+\tau^r.$$ 
Let  
$$
\cM=\left\{ a+b\tau^s\mid a, b\in \F_{q^r}\right\} \subset \End(\phi). 
$$
We further assume that $s<r$, so that $u=a+b\tau^s\in \cM$ is the expansion of $u$ with respect to the standard basis of $M_\phi$. In that case,   
one computes that the matrix $U$ by which $u$ acts on $M_\phi$ with respect to the standard basis is  
\begin{equation}\label{eqU2}
	U=\diag(a, a^q, \dots, a^{q^{r-1}})+\begin{bmatrix} & D_{s}\\ I_{r-s} & \end{bmatrix}\diag(b, b^q, \dots, b^{q^{r-1}}), 
\end{equation}
where 
$$
D_s =\diag(T-t, T-t^q, \cdots, T-t^{q^{s-1}})\in \Mat_{s}(k[T]). 
$$
Next, one computes that, when $r\neq 2s$, the determinant of $U$ in \eqref{eqU2} is equal to 
$$
\det U =\Nr_{\F_{q^r}/\F_q}(a) +(-1)^{r-1} \Nr_{\F_{q^r}/\F_q}(b)\cdot \Nr_{\F_{q^s}/\F_q}(T-t).  %(T-t)(T-t^q)\cdots (T-t^{q^{s-1}}). 
$$

We want to show that for an appropriate choice of $t$ and $\fp$, all $u\in \cM$ act invertibly on $\phi[\fp]$. First, note that given any prime 
$\fq\lhd A$ of degree $s$, we can choose $t\in \F_{q^s}$ so that $\fq=\Nr_{\F_{q^s}/\F_q}(T-t)$. Next, for a prime $\fp\lhd A$, by Dirichlet's theorem, any element of $\F_\fp$ is the residue modulo $\fp$ of some prime $\fq$. Thus, if $\deg(\fp)\geq 2$, 
there is always $\fq$ which is \textit{not} congruent modulo $\fp$ 
to an element of $\F_q$. Assume $\fp$ and $\fq$ are chosen with this property.  

If $u\in \cM$ is not invertible on $\phi[\fp]$, then $\fp$ divides
$$\det(U)= \Nr_{\F_{q^r}/\F_q}(a) +(-1)^{r-1} \Nr_{\F_{q^r}/\F_q}(b)\cdot \fq.$$
 If $b=0$, this is obviously not possible, so we assume $b\neq 0$. With this assumption, if $\fp$ 
divides $\det(U)$, then $\fq$ is congruent modulo $\fp$ to an element in $\F_q$, which leads to a contradiction. Thus, we can always choose $s$ and $t\in \F_{q^s}$ 
so that the images of nonzero elements of $\cM$ in $\End_{\F_\fp}(\phi[\fp])$ are invertible. Finally, note that $\dim_{\F_q}(\cM)=2r$, so 
by \eqref{eqSFdim} we get a semifield code when $\deg(\fp)=2$. 
\end{example}
%=======================================================

\section{Reinterpretation of Sheekey's construction}\label{secSheekey} Let $n\geq 1$ be an integer, let $k=\F_{q^n}$ be the degree 
$n$ extension of $\F_q$, and let $R=\twist{k}$ be the twisted polynomial ring with coefficients in $k$. 
The  center of $R$ is $A'\colonequals \F_q[\pi]$, where $\pi\colonequals \tau^n$. Note that $R$ is a free $A'$-algebra of rank $n^2$. 
%Given a polynomial $f$ in $A$, $A'$, or $R$, we denote by $f_0$ its constant coefficient, i.e., $f_0=f(0)$. 

\begin{defn} A \textit{central left multiple} of nonzero $u\in R$ is a polynomial $f\in A'$ such that $f=w u$ for some $w\in R$. 
	The \textit{minimal central left multiple} of $u$ is the unique central left multiple of $u$ which is monic and has minimal degree in $\pi$. 
\end{defn}

\begin{rem} The existence of a central left multiple can be seen, for example, by taking the norm of $u$ from $\F_q[u, \pi]$ into $\F_q[\pi]$. 
	The uniqueness of the minimal central left multiple follows from the minimality of degree and monic assumption, since the difference of two central left multiples is still a central left multiple.  Finally, using the division algorithm in $A'$, we see that any central left multiple of $u$ 
	is divisible by the minimal central left multiple. 
\end{rem}

Given $0\neq u\in R$, we define a Drinfeld module $\psi\colon A\to \twist{k}$ by $\psi_T=u$. Let $r_u\colonequals \deg_\tau u = \rank(\psi)$. 
Note that $\pi\in \End(\psi)$, since $\pi$ is in the center of $R$. 
Let $m_{u}(x)\in A[x]$ be the minimal polynomial of $\pi$ over $\F_q(u)$. 
Let $P_{u}(x)\in A[x]$ be the characteristic polynomial of $\pi$ acting on $T_\fp(\psi)$ for any $\fp\neq \chr_A(\psi)$. 
Let $\overline{m}_{u, T}(x)\in \F_q[x]$ be the polynomial obtained by reducing the coefficients of $m_{u}(x)$ modulo $T$, and 
define $\overline{P}_{u, T}(x)\in \F_q[x]$ similarly. Denote $L=\F_q(u)$, $\widetilde{L}=\F_q(u, \pi)$, and 
$$
d_u=[\widetilde{L}:L]. 
$$ 

%As we mentioned in Section \ref{secTM}, the polynomial $\overline{P}_{f, \pi}(x)$ is the characteristic polynomial of $\pi$ acting as an $\F_q$-linear transformation on $\psi[T]$. 


\begin{lem}\label{lemCLM}
	With previous notation, $\overline{m}_{u, T}(\pi)$ is a central left multiple of $u$. 
\end{lem} 
\begin{proof}
	By definitions, $f=\overline{m}_{u, T}(\pi)$ annihilates $\psi[T]$, i.e., the roots of $u(x)$ are roots of $f(x)$.   If we show 
	that $\Ht(f)\geq \Ht(u)$, then $f=wu$ follows from Lemma 3.1.16 in \cite{PapikianGTM}. If 
	$\chr_A(\psi)\neq T$ (equivalently, $\Ht(u)=0$), then the desired inequality is trivially true. 
	
	Now assume $\chr_A(\psi)=T$. We recall that (see \cite[Thm. 4.2.2]{PapikianGTM}) 
	\begin{equation}\label{eqPandM}
		P_u(x)=m_u(x)^{r_u/d_u}. 
		\end{equation} 
	Moreover, \cite[Thm. 4.2.13]{PapikianGTM} implies that 
	$\overline{P}_{u, T}(x)=x^{\Ht(u)}\cdot g(x)$, where $g(x)\in \F_q[x]$ and $g(0)\neq 0$. Thus, 
	$$
	\Ht(f)=\frac{n\cdot d_u}{r_u}\Ht(u). 
	$$
	On the other hand, by  \cite[(4.1.3)]{PapikianGTM}, 
	$$
	\frac{n\cdot d_u}{r_u} = [\widetilde{L}:\F_q(\pi)]\geq 1. 
	$$
	Combining these two formulas, we get $\Ht(f)\geq \Ht(u)$ as desired. 
\end{proof}

%In this section, from now on we assume that $\partial(u)\neq 0$, or equivalently $\chr_A(\psi)\neq T$. 
\begin{comment}
Let $f=f_0+f_1\tau+\cdots+f_s\tau^s\in R$ be a polynomial of degree $s$. Define a Drinfeld module $\psi\colon A\to \twist{k}$ of rank $s$ 
by $\psi_T=f$. Assume $f_0\neq 0$, or equivalently $\chr_A(\psi)\neq T$. Note that $\pi\in \End(\psi)$, since $\pi$ is in the center of $R$. Let $P_{f, \pi}(x)\in A[x]$ be the characteristic polynomial of $\pi$ acting on $T_\fp(\psi)$ for any $\fp\neq \chr_A(\psi)$. 
In particular, we can take $\fp=T$. Let $\overline{P}_{f, \pi}(x)\in \F_q[x]$ 
be the polynomial obtained by reducing the coefficients of $P_{f, \pi}$ modulo $T$. As we mentioned in Section \ref{secTM}, the polynomial 
$\overline{P}_{f, \pi}(x)$ is the characteristic polynomial of $\pi$ acting as an $\F_q$-linear transformation on $\psi[T]$. 

\begin{lem}\label{lem1.10}
	The minimal central left multiple of $f$ is $m_{f, \pi}(\pi)$, where 
	$m_{f, \pi}(x)\in \F_q[x]$ is the minimal polynomial 
	of $\pi$ acting on the $\F_q$-vector space $\psi[T]\cong \F_q^s$. 
\end{lem}
\begin{proof}
	By definition, $m_{f, \pi}(\pi)$ annihilates $\psi[T]$.  
	Thus, by \cite[Lem. 3.1.16]{PapikianGTM}, we have a decomposition $m_{f, \pi}(\pi)=h f$ for some $h\in R$. This implies that $m_{f, \pi}(\pi)$ is a 
	central left multiple of $f$. Conversely, suppose $g(\pi)\in A'$ is a central left multiple of $f$, so that $g(\pi)=h'f$ for some $h'\in R$. 
	Then $g(\pi)$ annihilates $\psi[T]\cong \F_q^s$, so $g(\pi)=0$ 
	as a linear operator on $\psi[T]$. This implies that $m_{f, \pi}(x)$ divides $g(x)$. We conclude that $m_{f, \pi}(\pi)$ is the minimal 
	central left multiple of $f$.  
\end{proof}
\end{comment}

\begin{lem}\label{prop1.8} Assume $\partial(u)\neq 0$. 
	The polynomial $u$ is irreducible in $R$ if and only if $\overline{P}_{u, T}(x)$ is irreducible in $\F_q[x]$. 
\end{lem}
\begin{proof}
	Note that $\psi[T]$ is the set of roots of the $\F_q$-linear polynomial $\psi_T(x)=u(x)$, which we assume is 
	separable.  Since $T\neq \chr_A(\psi)$, the polynomial $\overline{P}_{u, T}(x)$ is the characteristic polynomial of $\pi$ acting on $\psi[T]$. 
	Therefore, the polynomial $\overline{P}_{u, T}(x)$ is reducible in $\F_q[x]$ if and only if $\psi[T]$ has a $\pi$-invariant $\F_q$-subspace. The 
	elements of an $\F_q$-subspace  $W\subseteq \psi[T]$ give an $\F_q$-linear polynomial $w(x)=\prod_{\alpha\in W}(x-\alpha)\in \overline{k}[x]$. 
By \cite[Lem. 3.1.16]{PapikianGTM}, there is $g\in \twist{\overline{k}}$ such that $u=g w$. 
The subspace $W$ is $\pi$-invariant if and only if the coefficients of $w(x)$ are in $k$, 
so the decomposition $u=gw$ takes place in $\twist{k}$. Thus, $\overline{P}_{u, T}(x)$ is reducible in $\F_q[x]$ if and only if $u$ is reducible in $R$. 
\end{proof}

%\begin{cor}\label{cor1.9}
%	If $f$ is irreducible in $R$, then $\overline{P}_f(x)$ is irreducible in $\F_q[x]$. In particular, $P_f(x)$ is irreducible in $A[x]$. 
%\end{cor}


\begin{thm}\label{thm2.5} Assume $\partial(u)\neq 0$. 
	If $u$ is irreducible in $R$, then the minimal central left multiple of $u$ is $\overline{P}_{u, T}(\pi)$. 
\end{thm}
\begin{proof}
	By Lemma  \ref{prop1.8}, the assumption that $u$ is irreducible implies that 
	$\overline{P}_{u, T}(x)\in \F_q[x]$ is irreducible. In that case, by \eqref{eqPandM}, $\overline{m}_{u, T}(x)$ is also irreducible and 
	$\overline{m}_{u, T}(x)= \overline{P}_{u, T}(x)$. Now, by Lemma \ref{lemCLM}, $\overline{m}_{u, T}(\pi)$ is 
	a central left multiple of $u$. Using the irreducibly of $\overline{m}_{u, T}(x)$, we conclude that it is the 
	minimal central left multiple. 
\end{proof}

\begin{rem} Theorem \ref{thm2.5} is equivalent to Theorem 3 in \cite{Sheekey}. We note that there is an unnecessary assumption in \cite{Sheekey} that  
$u$ is monic, but also a missing assumption that $u$ 
	is irreducible in $R$. % (the proof uses Theorem 2 (iii), where $f$ is assumed to be irreducible).}  
\end{rem}

Let $u=u_0+u_1\tau+\cdots+u_s\tau^s \in \twist{k}$ and assume $u_0\neq 0$, $u_s\neq 0$. 
Let $M_\psi$ be the motive of $\psi$ defined by $\psi_T=u$, and let  $S_\psi\in \Mat_{s}(k[T])$ be the matrix defined in \eqref{eqSphi}. From the calculations in 
$\S$\ref{secMotive}, one deduces that $\pi\in \End(\psi)$ acts on $M_\psi$ by the matrix 
$$S_{u, \pi}=S_\psi S_\psi^{(q)}\cdots S_\psi^{(q^{n-1})}.$$ 
Let $\overline{S}_{\psi}$ be the matrix $S_{\psi}$ modulo $T$, and 
$\overline{S}_{u, \pi}=\overline{S}_\psi\overline{S}_\psi^{(q)}\cdots \overline{S}_\psi^{(q^{n-1})}$. 
Note that 
$\det\overline{S}_{\psi}=(-1)^{1+s}(-u_0/u_s)=(-1)^s(u_0/u_s)$, 
so 
$$
\det \overline{S}_{u, \pi}= (-1)^{sn}N(u_0/u_s),
$$
where 
$$
N(a)\colonequals \Nr_{k/\F_q}(a)=a^{1+q+q^2+\cdots+q^{n-1}}. 
$$ 
By \eqref{eqPu}, the characteristic polynomial of the Frobenius endomorphism of $\psi$ is 
$P_{u}(x)=\det(xI_r-S_{u, \pi})\in A[x].$
This implies that the determinant of $\pi$ acting on $\psi[T]$ is 
$(-1)^{sn}N(u_0/u_s)$.  On the other hand, the same determinant is equal to $(-1)^s \overline{P}_{u, T}(0)$. Thus, 
\begin{equation}\label{eq1}
	N(u_0/u_s)=(-1)^{s(n-1)}\overline{P}_{u, T}(0). 
\end{equation}

%Then $\overline{S}_{u, \pi}$ is the matrix by which $\pi$ acts on $\psi[T]$. %, and the characteristic polynomial $\overline{P}_f(x)=\det(xI_r-\overline{S}_{f, \pi})\in \F_q[x]$ 
%of this matrix is obtained by reducing the coefficients of $P_f(x)$ modulo $T$. 




Now let $\phi$ be the Drinfeld module of rank $n$ defined by $\phi_T=\pi$. Note that $\chr_A(\phi)=T$ and for any $a(T)\in A$, we have $\phi_a=a(\pi)$. 
For this Drinfeld module the endomorphism ring is as large as possible, $\End(\phi)=R$, since $\pi$ 
is in the center of $R$. Let $\fp\in  A$ be a monic irreducible polynomial of degree $s$. Let $\fp_0\colonequals \fp(0)$ be the constant term of $\fp$. 
Assume $\fp\neq T$. Then $\phi[\fp]\cong \F_\fp^n$ and we have a natural homomorphism 
$$
\iota_\fp\colon R=\End(\phi)\To \End_{\F_\fp}(\phi[\fp])\cong \Mat_n(\F_\fp). 
$$

The next theorem is equivalent to Theorem 7 in \cite{Sheekey} (for $k=1$ in the notation of \textit{loc. cit.}).  

\begin{thm}\label{thm7} 
	Fix $\eta\in k$ and define 
	$$
	\cM = \{u_0+u_1\tau\cdots+u_s\tau^s\in R\mid  u_s=\eta u_0\}. 
	$$
	If $N(\eta)\cdot (-1)^{s(n-1)}\fp_0\neq 1$, then the image of $\cM$ in $\End_{\F_\fp}(\phi[\fp])$ 
	is a semifield code. 
\end{thm}
\begin{proof}
	The dimension of $\cM$ over $\F_q$ is $ns$ (there is a dependency between the leading and the constant coefficients of the 
	polynomials in $\cM$). Therefore, by \eqref{eqSFdim}, to show that $\cM$ defines a semifield code 
	we need to show that every $0\neq u\in \cM$ acts invertibly on $\phi[\fp]$. 
	
	If $\rank(\iota_\fp(u))<n$, then the kernel of $u\colon \phi[\fp]\to \phi[\fp]$ is nonzero. If $0\neq \alpha\in \ker(u|_{\phi[\fp]})$, then $b\ast \alpha$ 
	is also in $\ker(u|_{\phi[\fp]})$ for all $b\in A$. Thus, $\ker(u)$ contains the 1-dimensional $\F_\fp$-subspace spanned by $\alpha$. But $\deg_\tau(u)\leq s$, 
	so this is possible only if $\deg_\tau(u)=s$ and $\phi_\fp=\fp(\pi)=wu$ for some $w\in R$. Moreover, we claim that $u$ must be irreducible in $R$. 
	If this is not the case, so $u=u_1u_2$ is a decomposition in $R$ with $\deg(u_i)<s$, then either $u_1$ or $u_2$ is not invertible on $\phi[\fp]$. But we just saw that 
	this is not possible. 
	
	Finally, if $u\in \cM$ is irreducible of degree $s$, then by Theorem \ref{thm2.5} the minimal central left multiple of $u$ is $\overline{P}_{u, T}(\pi)$, which is 
	irreducible in $A'$ of degree $s$. Thus, if $\fp(\pi)=wu$, then $\fp(\pi)=\overline{P}_{u, T}(\pi)$. But now, by \eqref{eq1}, we must have 
	$$
	N(1/\eta)=N(u_0/u_s)=(-1)^{s(n-1)}\fp_0,
	$$ 
	which contradicts the assumption of the theorem. 
\end{proof}

\begin{comment} 
There is a generalization of the fact that $\fp(x)=\overline{P}_{u, T}(x)$ used in the proof of the previous theorem. 
Let $\phi$ be a Drinfeld module over $k$ and $\fp\in A$ 
be a prime. Suppose $\phi_\fp= wu$, where $u\in \End(\phi)$ and $\partial u\neq 0$. If we define $\psi_T=u$, then $\phi_T\in \End(\psi)$. Let $\overline{P}_{\phi_T, u}(x)$ 
be the characteristic polynomial of the $\F_q$-linear operator $\phi_T$ acting $\psi[T]$. Assume $\phi_T$ acts irreducibly on $\psi[T]$, 
so that $\overline{P}_{\phi_T, u}(x)\in \F_q[x]$ is irreducible. (This is equivalent to assuming that no proper $\F_q$-subspace of $\ker(u)$ 
is the kernel of an isogeny from $\phi$ to some other Drinfeld module over an extension of $k$.) 
Then we claim that $\overline{P}_{\phi_T, u}(x)=\fp(x)$. To see this, note that 
$\overline{P}_{\phi_T, u}(\phi_T)$ annihilates $\psi[T]$ by the Cayley-Hamilton Theorem. 
On the other hand, $\phi_\fp= \fp(\phi_T)=wu$, so $\fp(\phi_T)$ also annihilates $\psi[T]$. By our assumption, $\overline{P}_{\phi_T, u}(x)$ 
is irreducible in $\F_q[x]$, so it is the minimal polynomial of $\phi_T$ acting on $\psi[T]$. Thus, $\overline{P}_{\phi_T, u}(x)$ divides $\fp(x)$ 
in $\F_q[x]$. Since both are irreducible and monic, we must have $\overline{P}_{\phi_T, u}(x)=\fp(x)$. 
\end{comment}


%----------------------------------------------------------

\section{A construction of semifield codes}\label{s2C}

In this section we generalize the construction in Example \ref{ssExample}. In this more general setting, explicit computations 
of determinants of endomorphisms become unwieldy, so we use various ad hoc arguments specific to the given situation to show the invertibility 
of endomorphisms acting on torsion points of Drinfeld modules.  

\subsection{Division algebras} 

In addition to the notation used in Section \ref{secSheekey}, let $\Delta=k(\tau)$ be the algebra over $K=\F_q(\pi)$ given in terms 
of generators and relations as follows:
\begin{align*}
	\Delta=k(\pi)\oplus k(\pi)\tau\oplus \cdots\oplus k(\pi)\tau^{n-1}, \\ 
	\tau^n=\pi, \quad \tau\alpha=\alpha^q\tau \quad \text{for all}\quad \alpha\in k. 
\end{align*} 
It is known that $\Delta$ is a central division algebra over $K$ of dimension $n^2$ and the local invariants of $\Delta$ are
zero at all places of $K$ except at $\pi$ and $1/\pi$, where the invariants are $1/n$ and $-1/n$, respectively; see \cite[Section 4.1]{PapikianGTM}. 

\begin{defn}
Let $\ell$  and $s$ be integers such that $n=\lcm(\ell, s)$. Let $g=\gcd(\ell,s)$. Denote 
$$
\Delta(\ell, s)=\F_{q^\ell}(\tau^s), \qquad \Delta(s, \ell)=\F_{q^s}(\tau^\ell), 
$$
so that $\Delta(n, 1)=\Delta$ and $\Delta(1, n)=\F_q(\pi)=K$. 
We consider $\Delta(\ell, s)$ and $\Delta(s, \ell)$ as division subalgebras of $\Delta$. 
\end{defn}

\begin{lem} Denote $K_g=\F_{q^g}(\pi)$. 
	\begin{enumerate}
		\item $Z(\Delta(\ell, s))=Z(\Delta(s,\ell))=K_g$, where $Z(\cdot)$ denotes the center of the corresponding algebra. 
		\item $\Cent_\Delta(\Delta(\ell, s))=\Delta(s, \ell)$ and $\Cent_\Delta(\Delta(s,  \ell))=\Delta(\ell, s)$, where $\Cent_\Delta(\cdot)$ denotes the centralizer in $\Delta$ 
		of the corresponding algebra. 
		\item $\dim_{K_g} \Delta(\ell, s) =(\ell/g)^2$ and $\dim_{K_g} \Delta(s, \ell) =(s/g)^2$. 
	\end{enumerate}
\end{lem}
\begin{proof} An element $\alpha=\sum_{j=0}^m a_j \tau^{j s}\in \twistt{\F_{q^\ell}}{\tau^s}$ commutes with $\tau^s$ if and only if all $a_j$ are in $\F_{q^s}\cap \F_{q^\ell}=\F_{q^g}$. Moreover, $\alpha$ commutes with all elements of $\F_{q^\ell}$ if and only if 
	$$a_j\neq 0\quad  \Longrightarrow \quad \ell \mid js. 
	$$
	Since $\ell \mid js$ is equivalent to $n\mid js$, we conclude that $Z(\Delta(\ell, s))\subseteq K_g$. The reverse inclusion is clear. This proves (1). 
	
	We have $s\cdot \ell = n\cdot g$, so 
	$$\pi=\tau^n = (\tau^s)^{\ell/g}. 
	$$
	Thus $\tau^s=\pi^{g/\ell}$ and, as a $K$-algebra, $\Delta(\ell, s)$ contains two linearly disjoint field extensions $\F_{q^\ell}K$ and $K(\pi^{g/\ell})$ of $K$ 
	of degrees $\ell$ and $\ell/g$, respectively. Thus, 
	$$
	\dim_K \Delta(\ell, s)\geq \ell^2/g. 
	$$
	Similarly, $\dim_K \Delta(s, \ell)\geq s^2/g$. It is easy to see that $\Delta(s, \ell)\subseteq \Cent_\Delta(\Delta(\ell, s))$. By the Double Centralizer Theorem, 
	$$
	\dim_K \Delta = \dim_K \Delta(\ell, s)\cdot \dim_K 	\Cent_\Delta(\Delta(\ell, s)). 
	$$
	Thus,
	$$
	n^2\geq ( \ell^2/g) ( s^2/g) = (\ell\cdot s/g)^2=n^2,
	$$
	so equalities must hold throughout. This proves (2) and (3). 
\end{proof}

\begin{prop}\label{prop3.2} Let $0\neq u\in \Delta(s, \ell)$, let $L=K(u)$, and let $D(u)=\Cent_\Delta(u)$. 
	\begin{enumerate}
		\item We have 
		$$
		\dim_L D(u) \geq (\ell/g)^2. 
		$$
		\item There is $u\in \Delta(s, \ell)$ such that $\dim_L D(u)=(\ell /g)^2$. Moreover, $$D(u)=L\otimes_{K_g} \Delta(\ell,s).$$ 
	\end{enumerate}
	%There is a Drinfeld module $\psi\colon A\to \twistt{\F_{q^s}}{\tau^\ell} \subseteq \twistt{\F_{q^n}}{\tau}$ such that $\End_k^\circ(\psi)$ is generated over $\F_q(T, \pi)$ by $\langle \F_{q^\ell}, \tau^s\rangle$. 
\end{prop}
\begin{proof} Because $K$ is the center of $\Delta$, we have $D(u)=\Cent_\Delta(L)$, so by the Double Centralizer Theorem (cf. 
	\cite[Thm. 1.7.22]{PapikianGTM}) $D(u)$ is a central division algebra over $L$ with 
	$$
	n^2=[D(u):K]\cdot [L:K]. 
	$$
	Hence $\dim_L D(u)=(n/[L:K])^2$. 
	
	On the other hand, we have the inclusions $K\subseteq K_g \subset \Delta(s, \ell)$. Moreover, $[K_g:K]=g$ and $[\Delta(s, \ell):K_g]=(s/g)^2$. 
	We proved that $K_g$ is the center of $\Delta(s, \ell)$. Hence a maximal subfield $M$ of $\Delta(s, \ell)$ has degree $s/g$ over $K_g$, and  
	$[M:K]=s$. In particular, $[L:K]\leq s$. From this inequality we get 
	$$\dim_L D(u)\geq (n/s)^2=(\ell/g)^2.$$
	This proves (1). 
	
	To prove (2), fix any maximal subfield $L\subset \Delta(s, \ell)$. Since $K$ is a function field of transcendence degree $1$ over a finite field, the Primitive Element Theorem (cf. \cite[Thm. 1.5.19]{PapikianGTM}) implies the existence of $u$ such that $L=K(u)$. For this $u$, $[L:K]=s$, so $\dim_L D(u)=(\ell/g)^2$. 
	On the other hand, we obviously have 
	$$
	L\otimes_{K_g}\Delta(\ell, s)\subseteq D(u). 
	$$
	Since the dimension of $\Delta(\ell, s)$ over $K_g$ is $(\ell/g)^2$, we have $\dim_L L\otimes_{K_g}\Delta(\ell, s) = (\ell/g)^2$. Comparing the 
	dimensions  we conclude that $L\otimes_{K_g}\Delta(\ell, s)=D(u)$. 
\end{proof}

\begin{rem}
	Note that in Proposition \ref{prop3.2} we can always scale $u$ by an element of $K$ to make it integral, i.e., to lie in $\twistt{\F_{q^s}}{\tau^\ell}$, 
	without affecting the claims about the dimension of $D(u)$. %The above proposition proves Lemma 3 and answers the question in Remark 4. 
\end{rem}

%More generally, we can take $D$ to be a subalgebra of $\Delta$ and $D'=\Cent_\Delta(D)$. The previous arguments apply in this setup, with 
%$K'=Z(D)=Z(D')=D\cap D'$ playing the role of $K_g$. Now take $\phi\colon A\to D$ and the message space to be $D'$. I am not sure if this leads to 
%anything useful. 

\begin{comment}
	%	\textcolor{blue}{What are the $K$-subalgebras of $\Delta$?} 
	Suppose $\phi\colon A\to \twist{k}$ is a Drinfeld module. 
	By definition, 
	$$
	D\colonequals \End(\phi)\otimes_A F=\Cent_\Delta(\phi_T) = \Cent_\Delta(\F_q(\pi, \phi_T)). 
	$$
	Hence $Z(D)=\F_q(\pi, \phi_T)$ and $D$ is the centralizer in $\Delta$ of its center, so in some sense it is 
	``maximal". (The subalgebra $\Delta(\ell, s)$ does not have this property.) 
	
	Suppose conversely that $D$ is a subalgebra of $\Delta$  
	such that $D=\Cent_\Delta(L)$, where $L=Z(D)$. Then $L$ is a finite field extension of $K$, and $D$ is a central division algebra 
	over $L$. By the Primitive Element Theorem, we can choose $u$ such that $L=K(u)$. After scaling, we may assume $u\in \twist{k}$. 
	Now $$D=\Cent_\Delta(L)= \Cent_\Delta(K(u))=\Cent_\Delta(u).$$
	If we define a Drinfeld module $\phi$ by $\phi_T=u$, then $\End(\phi)\otimes_A F=D$. Thus, $D$ is the endomorphism algebra 
	of a Drinfeld module. Note that in this construction $\chr_A(k)$ is the kernel of the $\F_q$-algebra homomorphism $A\to k$ defined by 
	$T\mapsto u_0$, so $\chr_A(k)$ \textit{depends} on $D$. We also note that we don't have much control over the rank of $\phi$. 
\end{comment}


%======================================================

\subsection{Key technical lemma}\label{sKeyLemma}

Let $\phi\colon A\to \twist{k}$ be a Drinfeld module. 
%Let $u\in \End(\phi)$ be an endomorphism with $\deg_\tau u>0$.  Let $\psi\colon \F_q[T]\to \twist{k}$ be the Drinfeld module defined by $\psi_T=u$. Assume $u_0\neq 0$ (or equivalently, $\chr_A(\psi)\neq T$). Let $m_{u, T}(x)\in \F_q[x]$ be the minimal polynomial of $\pi$ acting on the $\F_q$-vector space $\psi[T]\cong \F_q^{\deg_\tau(u)}$. By Lemma \ref{lem1.10}, $m_{u, T}(\pi)\in \F_q[\pi]$ is the minimal central left multiple of $u\in R$. 
Let $\fp\in A$ be an irreducible monic polynomial not equal to $\chr_A(\phi)$. 
Let $m_{\phi, \fp}(x)\in \F_\fp[x]$ be the minimal polynomial of $\pi$ acting on the $\F_\fp$-vector space 
$\phi[\fp]\cong \F_\fp^{\deg_\tau(\phi_T)}$. We take the norm 
$$
\Nr_{\F_\fp/\F_q}(m_{\phi, \fp}(x))= m_{\phi, \fp}(x)\cdot m_{\phi, \fp}^{(q)}(x)\cdots m_{\phi, \fp}^{(q^{\deg_T(\fp)-1})}(x)  
$$
to obtain a polynomial in $\F_q[x]$. Recall that given $0\neq u\in \End(\phi)$, we denoted 
$L=\F_q(u)$, $\widetilde{L}=\F_q(u, \pi)$, and $d_u=[\widetilde{L}:L]$.  

\begin{lem}\label{lemPrelim1}  If all 
	irreducible factors of 	$\Nr_{\F_\fp/\F_q}(m_{\phi, \fp}(x))$ in $\F_q[x]$ have degrees $> d_u$, then
	$u$ acts invertibly on $\phi[\fp]$. 
\end{lem} 

\begin{proof} We argue by contradiction. Suppose there is $0\neq \alpha\in \phi[\fp]$ is such that $u(\alpha)=0$. 
	By Lemma \ref{lemCLM}, we have a decomposition $\overline{m}_{u, T}(\pi)=wu$ in $R$, and $d_u=\deg_x\overline{m}_{u, T}(x)$ 
	from definitions. Thus, we have $\overline{m}_{u, T}(\pi)(\alpha)=0$.  This implies that the $\F_q$-span 
	of $\alpha, \pi(\alpha), \dots, \pi^{d_u-1}(\alpha)$ is a $\pi$-invariant subspace of $\oF_q$ of dimension $\leq d_u$. 
	
	On the other hand, the minimal polynomial of $\pi$ acting on $\phi[\fp]$,  
	as an $\F_q$-vector space, is $\Nr_{\F_\fp/\F_q}(m_{\phi, \fp}(x))$. The previous paragraph implies that 
	$\Nr_{\F_\fp/\F_q}(m_{\phi, \fp}(x))$ 
	has an irreducible factor of degree $\leq d_u$ in $\F_q[x]$, which contradicts the assumption of the lemma. 
\end{proof}

Let $m_\phi(x)$ be the minimal polynomial of $\pi\in \End(\phi)$ over $F=\F_q(\phi_T)$. Then $m_\phi(x)$ 
is a monic irreducible polynomial in $A[x]$ of degree $d_\phi\colonequals [\F_q(\pi, \phi_T):\F_q(\phi_T)]$. 
Let $\overline{m}_{\phi, \fp}(x)\in \F_\fp[x]$ be the polynomial obtained from $m_\phi(x)$ by reducing its coefficients 
modulo $\fp$. 


\begin{lem}\label{cor2.2} 
	Assume  $\overline{m}_{\phi, \fp}(x)$ is irreducible in $\F_\fp[x]$ and $\overline{m}_{\phi, \fp}(0)$ generates $\F_\fp$ over $\F_q$. 
If $d_u< d_\phi\cdot \deg_T(\fp)$,  
	then $u$ acts invertibly on $\phi[\fp]$. 
\end{lem}
\begin{proof} Note that $\overline{m}_{\phi, \fp}(\pi)=0$ on $\phi[\fp]$, so $m_{\phi, \fp}(x)$ divides $\overline{m}_{\phi, \fp}(x)$. 
	The irreducibility of this latter polynomial then implies that $m_{\phi, \fp}(x)=\overline{m}_{\phi, \fp}(x)$. 
	Hence 
	$$
	\deg_x\Nr_{\F_\fp/\F_q}(m_{\phi, \fp}(x)) = \deg_x(\overline{m}_{\phi, \fp}(x)) \cdot \deg_T(\fp)=d_\phi\cdot \deg_T(\fp). 
	$$
	
	We claim that $\Nr_{\F_\fp/\F_q}(m_{\phi, \fp}(x))$ is irreducible over $\F_q$; assuming this, the invertibility of $u$ 
	on $\phi[\fp]$ follows from Lemma \ref{lemPrelim1}.  
	Let $g(x)$ be a monic irreducible divisor of $\Nr_{\F_\fp/\F_q}(\overline{m}_{\phi, \fp}(x))$ in $\F_q[x]$. Considering $g(x)$ as a polynomial in $\F_\fp[x]$, 
	the irreducibility of all $\overline{m}_{\phi, \fp}^{(q^i)}(x)$ implies that $g(x)=\overline{m}_{\phi, \fp}^{(q^i)}(x)$ for some $0\leq i\leq \deg(\fp)-1$. 
	But then $\overline{m}_{\phi, \fp}(0)^{q^i}$ lies in $\F_q$,  which contradicts the assumption that $\overline{m}_{\phi, \fp}(0)$ generates $\F_\fp$ over $\F_q$. 
\end{proof}

%=============================================

\subsection{The construction of semifield codes} Recall that $\ell$ and $s$ are integers such that $n=\lcm(\ell, s)$, and we denoted 
$g=\gcd(\ell, s)$. 

Define $\phi\colon A\to \twist{k}$ by 
\begin{equation}\label{eqDefphi}
\phi_T=\sum_{i=0}^r a_i \tau^{\ell i}, \qquad a_i\in \F_{q^s}, 
\end{equation}
so that the image of $\phi$ lies in $\Delta(s, \ell)$ and $\phi$ has rank $r\ell$. 

Set 
$$
\cM=\left\{\sum_{i=0}^e b_i\tau^{s i}\mid b_i\in \F_{q^\ell}\right\}\subset \Delta(\ell, s)\subset \End(\phi). 
$$
Let $\fp\lhd A$ be a prime different from $\chr_A(\phi)$ and let $d\colonequals \deg(\fp)$. 
We want to give conditions on the prime $\fp$ and the Drinfeld module $\phi$ such that 
the image of $\cM$ in $\End_{\F_\fp}(\phi[\fp])\cong  \Mat_{r\ell}(\F_\fp)$ is a semifield code. 

Since $$\dim_{\F_q}\cM=(e+1)\ell, $$ to get a semifield code from $\cM$ we first of all need 
$
r\ell d= (e+1)\ell$; see \eqref{eqSFdim}. Thus,  
$$
\boxed{e+1=r d}
$$


\begin{lem}\label{lem1.10M} If $0\neq u\in \cM$, then $d_u\leq eg$. 
\end{lem}
\begin{proof} 
	Define a Drinfeld module $\psi$ over $k$ by setting $\psi_T=u$. Applying Theorem 4.1.3 in \cite{PapikianGTM} to $\psi$, we deduce that  
	$\widetilde{L}$ is the center of $D=\Cent_\Delta(u)$ and 
	$$
	\dim_{\widetilde{L}}(D)=\left(\frac{\deg_\tau(u)}{d_u}\right)^2 \leq \left(\frac{se}{d_u}\right)^2. 
	$$
	On the other hand, since $u\in \Delta(\ell, s)$, Proposition \ref{prop3.2} implies that 
	$
	\dim_{\widetilde{L}}(D) \geq (s/g)^2$. 
	Thus, $d_u\leq eg$. 
\end{proof}

\begin{assumptions}
Let $\tF=\F_q(\phi_T, \pi)$ and $F=\F_q(\phi_T)$. Suppose 
\begin{enumerate}
	\item $\overline{m}_{\phi, \fp}(x)$ is irreducible in $\F_\fp[x]$.
	\item $\overline{m}_{\phi, \fp}(0)$ generates $\F_\fp$ over $\F_q$. 
	\item $eg<d\cdot [\tF:F]$. 
\end{enumerate}
Then, by  Lemma \ref{cor2.2} and Lemma \ref{lem1.10M}, any $0\neq u\in \cM$ acts invertibly on $\phi[\fp]$. 
\end{assumptions}

First, we examine the implication of assumption (3). Applying Lemma \ref{lem1.10M} to $\phi$ leads to 
\begin{equation}\label{eqtF1}
	[\tF:F]\leq rg. 
\end{equation}
On the other hand, from $eg<d[\tF:F]$ and $e+1=rd$, we get 
\begin{equation}\label{eqtF2}
	eg<d[\tF:F] = \frac{(e+1)}{r}[\tF:F]. 
\end{equation}
Combining \eqref{eqtF1} and  \eqref{eqtF2}, we get 
$$
\frac{e}{e+1} rg < [\tF:F]\leq rg. 
$$
We want to allow $d$ (and thus $e$) to be arbitrarily large, so we are forced to assume that 
$$
\boxed{ [\tF:F]=rg}
$$
By \eqref{eqtF1}, this is the maximal possible degree for the extension $\tF/F$. On the other hand, by 
Proposition \ref{prop3.2}, one can choose $\phi_T\in \Delta(s, \ell)$,  
so that this equality holds. We assume that $\phi$ is chosen to have this property. 

\begin{rem} Let $\fl=\chr_A(\phi)$. 
	The \textit{height} of $\phi$ is defined as $H(\phi)=\Ht(\phi_\fl)/\deg(\fl)$. By \cite[Lem. 3.2.11]{PapikianGTM}, $H(\phi)$ 
	is a positive integer. Now note that $\phi$ 
	defined by \eqref{eqDefphi} can be considered as a $\F_{q^\ell}[T]$-Drinfeld module $\phi'$ of rank $r$. 
	Since $H(\phi)=\ell H(\phi')$, we conclude that $\ell\leq H(\phi)$. On the the other hand, by \cite[Prop.  4.1.8]{PapikianGTM}, 
	$H(\phi)\geq r\ell/[\tF:F]$. Thus, if $H(\phi)=\ell$, then $r\leq [\tF:F]\leq rg$. This means that when $g=1$, without appealing to Proposition \ref{prop3.2}, 
	we can choose $\phi$ for which $[\tF:F]=r$ is maximal by choosing $\phi'$ to be \textit{ordinary}, i.e., having $H(\phi')=1$.  Since ``most" Drinfeld modules over finite fields are ordinary, such a choice is always possible. 
\end{rem}

It remains to show that there are primes $\fp$ for which assumptions (1) and (2) hold. 

\begin{lem} Assume $g=1$ and $r$ is not divisible by the characteristic of $\F_q$. In addition, assume that either $r$ is prime, or 
that $r$ is coprime to $s$. Then the Galois group of $m_\phi(x)$ contains a cycle of maximal length $r$. 
\end{lem}
\begin{proof}
	Let $M$ be the splitting field of $m_\phi(x)$. Since $r=\deg m_\phi(x)$ is not divisible by the characteristic of $F$, the extension $M/F$ is Galois. 
	We claim that $\Gal(M/F)$, as a subgroup of the permutation group $S_r$ of the roots of $m_\phi(x)$, contains a cycle of length $r$. 
	We will show that the decomposition subgroup $G_\infty\subseteq \Gal(M/F)$ at the place $\infty=1/T$ contains such a cycle. 
	
	By \cite[Thm. 4.1.3]{PapikianGTM}, $m_\phi(x)$ remains irreducible over the completion $\Fi$ of $F$. Let 
	$M'$ be the splitting field of $m_\phi(x)$ over $\Fi$. Then $G_\infty\cong \Gal(M'/\Fi)$. 
	Let $\tinf$ be the unique place of $\tF$ over $\infty$, and let $\tF_{\tinf}$ be the completion of $\tF$ at $\tinf$. We have 
	$r=[\tF_{\tinf}:\Fi]=e(\tF_{\tinf}/\Fi) f(\tF_{\tinf}/\Fi)$, where $e(\tF_{\tinf}/\Fi) $ is the ramification index and $f(\tF_{\tinf}/\Fi)$ 
	is the residual degree. 
	
	If $r$ is a prime, then either  $e(\tF_{\tinf}/\Fi)=r$ or $f(\tF_{\tinf}/\Fi)=r$. In the second case, $M'=\tF_{\tinf}$ is an unramified extension of degree $r$, so $\Gal(M'/\Fi)\cong \Z/r\Z$. In the first case, $\tF_{\tinf}/\Fi$ is a totally tamely ramified extension. 
	If $r$ is not necessarily prime, we note that the  normalized 
	valuations of the roots of $m_\phi(x)$ at $\infty$ are equal to $-n/r\ell=-s/r$; see \cite[Thm. 4.2.7]{PapikianGTM}. 
	Therefore, if $r$ is coprime to $s$, then again this implies that $\tF_{\tinf}/\Fi$ is a totally tamely ramified extension. 
	
	We are reducing to showing that if $\tF_{\tinf}/\Fi$ is a totally tamely ramified extension of degree $r$, 
	then $\Gal(M'/\Fi)$ contains a cycle of length $r$. By a well-known structure theorem for totally tamely ramified extensions 
	of local fields (cf. \cite[Prop. 2.6.7]{PapikianGTM}), $M'$ is the splitting field of $x^r-\varpi_\infty$, where $\varpi_\infty$ 
	is a uniformizer of $\Fi$. Therefore, $M'$ contains a primitive $r$-th root $\zeta_r$ of $1$, and the automorpism $M'\to M'$ defined by 
	$\varpi_\infty^{1/r}\mapsto \zeta_r\varpi_\infty^{1/r}$ is a cycle of length $r$. 
\end{proof}

Assume the Galois group of $m_\phi(x)$ contains a cycle of length $rg$ (this is always the case under the assumptions on $rg$ of the previous lemma). 
By the effective Chebotarev density theorem (cf. \cite{MuSch}), 
the number of primes in $A$ of degree $N\gg 0$, which are unramified in the Galois closure of $\tF/F$ and whose corresponding Frobenius is in the conjugacy class of a maximal cycle in $\Gal(m_\phi(x))$, is $\geq cq^N/N$ for some nonzero constant $c$ (not depending on $N$). Denote the set of such primes by $\cP_N$. 

For $\fp\in \cP_N$ the reduction of $m_\phi(x)$ modulo $\fp$ is irreducible (i.e., (1) holds), so we now concentrate on condition (2).  
Let $\fq=\chr_A(\phi)$. By \cite[Thm. 4.2.2]{PapikianGTM}, $P_\phi(x)=m_\phi(x)^{r\ell/rg}=m_\phi(x)^{\ell/g}$. Moreover, by 
\cite[Thm. 4.2.7]{PapikianGTM}, up to an $\F_q^\times$-multiple, $P_\phi(0)$ is equal to $\fq^{n/\deg(\fq)}$. Hence, up to an $\F_q^\times$-multiple, 
$m_\phi(0)$ is equal to $\fq^{s/\deg(\fq)}$. %(In particular, $\deg(\fq)$ necessarily divides $s$ -- a restriction that is not obvious from the outset.)

Denote $f(T)=\fq^{s/\deg(\fq)}\in A$. Note that $\deg_T f=s$. Choose a root of each $\fp\in \cP_N$ in $\overline{\F}_q$ and call the set of these roots $\cR_N$. We have 
$\# \cR_N\geq  cq^N/N$. We want to show that $f(T)$ modulo $\fp$ generates $\F_\fp=\F_{q^N}$ over $\F_q$ for at least one $\fp\in \cP_N$. 
This is equivalent to requiring that for some $\alpha\in \cR_N$, the value $f(\alpha)$ 
does not belong to a proper subfield of $\F_{q^N}$.
The image of the map $\cR_N\to \F_{q^N}$, $\alpha\mapsto f(\alpha)$, has size at least $\# \cR_N/\deg(f)$ since the preimages of this map for $\beta\in \F_{q^N}$ 
are the roots of $f(x)-\beta$, and there at most $\deg(f)$ such roots. Finally, note that the union of proper subfields of $\F_{q^N}$ has size at most $O(q^{N/2})$. 
Since 
$$
c\frac{q^N}{N\cdot s} > O(q^{N/2})
$$ 
when $N$ is sufficiently large, we see that we can always find $\fp$ such that the properties (1) and (2) are satisfied. 

 \subsection{The nuclear parameters of the code} As in the previous subsection, let  $r\geq 1$, $d=\deg \fp$, 
\[
\phi_T=\sum_{i=0}^r a_i \tau^{\ell i} \in \twistt{\F_{q^s}}{\tau^\ell},
\]
and 
\begin{equation}\label{eqcM}
\cM=\left\{\sum_{i=0}^{rd-1} b_i\tau^{s i}\mid b_i\in \F_{q^\ell}\right\}. 
\end{equation}
Assume all nonzero elements of $\cM$ act invertibly on $\phi[\fp]$, so that $\cM$ gives a semifield code in $\End_{\F_\fp}(\phi[\fp])$. 

\begin{defn}\label{defNP}
Define \textit{left (right) idealizer, centralizer, and center} of $\cM$ as follows:  
\begin{align*}
	I_l(\cM) &=\{f\in \End_{\F_\fp}(\phi[\fp])\mid f\cM\subseteq \cM\},\\ 
	I_r(\cM) &=\{f\in \End_{\F_\fp}(\phi[\fp])\mid \cM f\subseteq \cM\},\\ 
	C(\cM) &=\{f\in \End_{\F_\fp}(\phi[\fp])\mid mf=fm \text{ for all }m\in \cM\},\\ 
	Z(\cM) &=I_l(\cM)\cap C(\cM). 
\end{align*}
By \cite[p. 436]{Sheekey}, each of these sets is a field extension of $\F_q$. 
The \textit{nuclear parameters} of $\cM$ is the tuple
\[
(\dim_{\F_q}\cM, \dim_{\F_q} 	I_l(\cM),  \dim_{\F_q} 	I_r(\cM),  \dim_{\F_q} 	C(\cM),  \dim_{\F_q} 	Z(\cM)). 
\]
By \cite[Prop. 4]{Sheekey}, equivalent codes have the same nuclear parameters, so this tuple gives an 
equivalence invariant of the code. 
\end{defn}


The ``expected" parameters of $\cM$ are 
\[
\boxed{(rd\ell, \ell, \ell, rdg, g)}
\]
More precisely, it is clear that $\F_{q^\ell}\subseteq I_l(\cM)$, $\F_{q^\ell}\subseteq I_r(\cM)$, and $\F_{q^g}\subseteq Z(\cM)$. 
The centralizer contains the image of $\F_{q}[\phi_T, \pi]$ in $\End_{\F_\fp}(\phi[\fp])$. 
Recall that we are assuming that $\deg m_\phi(x)=rg$ and $m_\phi(x)$ 
is irreducible modulo $\fp$. Under these assumptions, we have 
\[\F_{q}[T, \pi]/\fp\cong \F_\fp[x]/(\overline{m}_{\phi, \fp}(x))\cong \F_{\fp^{rg}}\cong \F_{q^{drg}}. 
\]
(Note that the image of $\F_{q}[\phi_T, \pi]$ in $\End_{\F_\fp}(\phi[\fp])$ already contains  $\F_{q^g}\subseteq Z(\cM)\subseteq C(\cM)$.) 
The expectation is that there are no sporadic elements in the idealizers and the centralizer, which in principle might occur because we are working modulo $\fp$. In any case, we have the following: 


\begin{lem} If $s< \ell$, then 
\[ 
	I_l(\cM)=I_r(\cM)=\F_{q^\ell}, \qquad Z(\cM)=\F_{q^g}. 
\]
\end{lem}
\begin{proof} Since $1\in \cM$, if $f\in I_l(\cM)$, then we can assume that $f\in \cM$. To prove that 
	$I_l(\cM)=\F_{q^\ell}$, it is enough to prove that $f$ cannot have positive degree in $\tau$. Let $w=\deg_\tau(f)$. 
	Note that $w\leq (rd-1)s$ and $s\mid w$.  
	If $w>0$, then $\tau^{srd-w}\in \cM$ but $f \tau^{srd-w}\not\in \cM$, since $srd<\ell rd=\deg_\tau(\phi_\fp)$. This 
	leads to a contradiction. The argument for $I_r(\cM)$ is similar. 
	The elements of $\F_{q^\ell}$ that commute with all elements of $\cM$ are those in $\F_{q^g}$, so $Z(\cM)=\F_{q^g}$. 
\end{proof}

\begin{comment}
\begin{color}{red}Argument for the centralizer. Suppose that the centralizer has dimension larger than $rdg$. Say for example that such $\mathbb F_q$ dimension is $rdgh$. Now, the degree $rd-1$ on the message space is exactly tuned so that the centralizer allows for too many zeroes as long as there is one. So, if the centralizer was larger by a factor $h$ so one could take a larger $e$, up to $drh$? But this denies the Singleton bound, as our code is already MRD? I am not sure, I just wanted to ask whether you thought about this kind of argument.
\end{color}
\end{comment}

\begin{lem}\label{lemCentralizer}
    Assume $m_\phi(x)$ is irreducible modulo $\fp$, $g=1$, and $\ell$ is prime. Then $C(\cM)\cong \F_{q^{rd}}$. 
\end{lem}
\begin{proof}
    We have already seen above that $\F_{q^{rd}}\subseteq C(\cM)$. On the other hand, $C(\cM)$ is a field acting on \[\phi[\fp]\cong \F_{\fp}^{r\ell}\cong \F_q^{rd\ell},\] 
    i.e., $\phi[\fp]$ is a vector space over $C(\cM)$, so $C(\cM)\subseteq \F_{q^{rd\ell}}$. If $C(\cM)$ is strictly larger than 
    $\F_{q^{rd}}$ and $\ell$ is prime, then $C(\cM)\cong \F_{q^{rd\ell}}$. 
    %so $\phi[\fp]$ is 1-dimensional over $C(\cM)$. 
    Denote the centralizer of $C(\cM)$ in $\End_{\F_\fp}(\phi[\fp])$ by $C(C(\cM))$. By 
    the Double Centralizer Theorem applied to $D=\Mat_{r\ell}(\F_\fp)$ 
    (see \cite[Cor. 7.13]{Reiner}), we have 
    \[
    [D:\F_\fp] = [C(\cM):\F_\fp] \cdot [C(C(\cM)):\F_\fp]. 
    \]
    This implies that $[C(C(\cM)):\F_\fp]=r\ell$. On the other hand, we have $C(\cM)\subseteq C(C(\cM))$, so $C(\cM)= C(C(\cM))$. 
    Note that $\cM\subseteq C(C(\cM))$. Since $\dim \cM = rd\ell = \dim C(\cM)$, we conclude that the image of $\cM$ in $\End_{\F_\fp}(\phi[\fp])$ is isomorphic to $\F_{q^{rd\ell}}$; in particular, it is commutative. But $\tau^\ell$ and $\F_{q^s}$ do not commute modulo $\phi_\fp$, so we get a contradiction. Therefore, $C(\cM)\cong \F_{q^{rd}}$.
\end{proof}

\begin{example} The calculations in this example were done using  \texttt{Magma} \cite{MR1484478}. 
	
	Let $q=3$, $\ell=2$ and $s=3$, so $n=6$. Let $\alpha$ be a root of $x^3-x+1$, which is irreducible in $\F_q[x]$.  Define 
	$$
	\phi_T=\alpha+\alpha^2\tau^2+\tau^4. 
	$$
	The minimal polynomial of $\pi=\tau^6$ over $\F_q(\phi_T)$ is 
	$$
	x^2-Tx-(T^3-T+1). 
	$$
	This polynomial modulo $\fp=T-1$ is $x^2+x-1$, which is irreducible over $\F_q$.
	
	It is not hard to check that 
	$$
	\cM=\{b_0+b_1\tau^3\mid b_0, b_1\in \F_{q^2}\}
	$$
	gives a semifield code in $\End_{\F_{T-1}}(\phi[T-1])\cong \Mat_4(\F_q)$. (For this, one just needs to check that all nonzero $b_0+b_1\tau^3$ 
	are coprime to $\phi_{T-1}=(\alpha-1)+\alpha^2\tau^2+\tau^4$ in $\twist{k}$, which is easy to do using the division algorithm; cf. \cite[Thm. 3.1.13]{PapikianGTM}.) Next, one verifies that no element $f$ of $\cM$ of positive degree satisfies \[f\cM\pmod{\phi_{T-1}}\subseteq \cM.\]
	Thus, $I_l(\cM)=\F_{q^2}$, and similarly $I_r(\cM)=\F_{q^2}$. 
    By Lemma \ref{lemCentralizer}, $C(\cM)=\F_{q^2}$.
    \begin{comment}
	To determine the centralizer, one first computes that the minimal polynomial of $\pi=\tau^6$ over $\F_q(\phi_T)$ is 
	$$
	x^2-Tx-(T^3-T+1). 
	$$
	This polynomial modulo $T-1$ is $x^2+x-1$, which is irreducible over $\F_q$. Hence the image of $A[\pi]$ in $\End(\phi[T-1])$ 
	is isomorphic to $\F_{q^2}$. On the other hand, the image of $A[\pi]$ lies in $C(\cM)$. As we mentioned, $C(\cM)$ is a field extension of 
	$\F_q$. Since this field embeds into $\Mat_4(\F_q)$, it acts on $\F_q^4$, so $C(\cM)$ is either $\F_{q^2}$ or $\F_{q^4}$. 
	Suppose $C(\cM)=\F_{q^4}$. By the Double Centralizer Theorem, the centralizer of $\F_{q^4}$ in $\Mat_4(\F_q)$ is $\F_{q^4}$ 
	itself. Comparing the dimensions, we conclude that $\cM=\F_{q^4}$. This leads to a contradiction, since multiplication on $\cM$ modulo 
	$\phi_{T-1}$ is not commutative. Thus, $C(\cM)=\F_{q^2}$. 
    \end{comment}
	Finally, the elements of $I_l(\cM)$ that commute with all other elements of $\cM$ are the elements of $\F_q$. Thus, $Z(\cM)=\F_q$. We conclude that 
	the parameters of our semifield code are 
	\[
	(4,2,2,2,1). 
	\] 
\end{example}

 %---------------------------------------------------------
 
 \subsection{Centralizers of matrices over extension fields} The torsion module $\phi[\fp]$ is most naturally an $\F_\fp$-vector space, as 
 it reflects the $A$-module structure given by $\phi$. On the other hand, $\phi[\fp]$ can also be considered as an $\F_q$-vector space. Then, similar 
 to Definition \ref{defNP},  
 one can define the idealizers and the centralizer of $\cM$ in $\End_{\F_q}(\phi[\fp])$. In fact, if one follows Definition 2 in \cite{Sheekey}, 
 then this is how these invariants should be defined in our context. 
 
 Fix an embedding $\iota\colon \F_\fp\hookrightarrow \Mat_d(\F_q)$, for example the regular representation. This induces an embedding 
 \[
 \iota\colon \Mat_r(\F_\fp) \longhookrightarrow \Mat_{rd}(\F_q)
 \]
 by applying $\iota$ entry-wise (or equivalently, by viewing $\F_\fp^r$ as an $\F_q$-vector space of dimension $rd$). 
 For our message space $\cM$, we wish to compare $\cI_\ell(\iota(\cM))$, $\cI_r(\iota(\cM))$, and 
 $C(\iota(\cM))$ with $\cI_\ell(\cM)$, $\cI_r(\cM)$, and $C(\cM)$, respectively.  Because, $1\in \cM$, it is easy to see that the left and right 
 idealizers do not change. The situation with the centralizer is less clear, 
 so we first obtain a general criterion that can be applied to answer the question. 
 
 For a matrix $M\in \Mat_r(\F_\fp)$, we wish to compare:
 \begin{itemize}
 	\item $C_{\Mat_{rd}(\F_q)}(\iota(M))$, the centralizer of $\iota(M)$ in $\Mat_{rd}(\F_q)$, and 
 	\item $\iota\left(C_{\Mat_r(\F_\fp)}(M)\right)$, the image under $\iota$ of the centralizer of $M$ in $\Mat_r(\F_\fp)$.  
 \end{itemize}
 
 \begin{lem}\label{lemC(F_p)}
 	The image of $\iota$ satisfies 
 	\[
 	\iota(\Mat_r(\F_\fp)) = C_{\Mat_{rd}(\F_q)}(\iota(\F_\fp)). 
 	\]
 \end{lem}
 \begin{proof}
 	View $\F_\fp^r$ as an $\F_q$-vector space of dimension $rd$. An $\F_q$-linear endomorphism of $\F_\fp^r$ lies 
 	in the image of $\iota$ if and only if it is $\F_\fp$-linear. But $\F_\fp$-linearity means precisely that the endomorphism 
 	commutes with scalar multiplication by every element of $\F_\fp$, i.e., it commutes with $\iota(\alpha)$ for all $\alpha\in \F_\fp$. 
 \end{proof}
 
 \begin{prop}\label{propCM}
 	Let $M\in \Mat_r(\F_\fp)$. The following are equivalent: 
 	\begin{itemize}
 		\item[(i)] $C_{\Mat_{rd}(\F_q)}(\iota(M))=\iota\left(C_{\Mat_r(\F_\fp)}(M)\right)$.
 		\item[(ii)]  $\iota(\F_\fp)\subseteq \F_q[\iota(M)]$. 
 	\end{itemize}
 \end{prop}
 \begin{proof} To simplify the notation, we will denote 
 	\[C(\iota(M))= C_{\Mat_{rd}(\F_q)}(\iota(M))\quad \text{and}\quad \iota(C(M))=\iota\left(C_{\Mat_r(\F_\fp)}(M)\right).
 	\] 
 	
 	If $N\in \Mat_r(\F_\fp)$ satisfies $NM=MN$, then applying $\iota$ (which is a ring homomorphism) gives $\iota(N)\iota(M)=\iota(M)\iota(N)$. 
 	Thus, the inclusion  
 	\[
 	\iota\left(C(M)\right) \subseteq C(\iota(M)) 
 	\]
 	holds unconditionally. 
 	
 	For the reverse inclusion, by Lemma \ref{lemC(F_p)}, we have 
 	\[
 	\iota\left(C(M)\right)  = 	\iota(\Mat_r(\F_\fp))  \cap C(\iota(M)) = C(\iota(\F_\fp)) \cap C(\iota(M)). 
 	\]
 	Therefore, the equality $\iota\left(C(M)\right)  = C(\iota(M))$ holds if and only if 
 	\[
 	C(\iota(M))\subseteq C(\iota(\F_\fp)). 
 	\]
 	
 	(ii)$\Rightarrow$(i): Suppose $\iota(\F_\fp)\subseteq \F_q[\iota(M)]$. If $X\in C(\iota(M))$, then $X$ commutes with every 
 	polynomial in $\iota(M)$, hence $X$ commutes with every element of $\F_q[\iota(M)]$. In particular, $X$ commutes 
 	with every element of $\iota(\F_\fp)$, so $X\in C(\iota(\F_\fp))$. 
 	
 	(i)$\Rightarrow$(ii):  Suppose $C(\iota(M))\subseteq C(\iota(\F_\fp))$. Taking centralizers reverses inclusions, so 
 	\[
 	C(C(\iota(\F_\fp)))\subseteq C(C(\iota(M))). 
 	\]
 	On the other hand, by the Double Centralizer Theorem, 
 	\begin{align*}
 	C(C(\iota(\F_\fp))) &= \iota(\F_\fp), \\ 
 	C(C(\iota(M))) &= \F_q[\iota(M)].  
 	\end{align*}
 	Therefore, $\iota(\F_\fp) \subseteq \F_q[\iota(M)]$. 
 \end{proof}
 
 We now give a more explicit characteriszation of condition (ii) in Proposition \ref{propCM}. 

\begin{prop}\label{propCM2}
	Let $M\in \Mat_r(\F_\fp)$. The following are equivalent:
	\begin{itemize}
		\item[(ii)]  $\iota(\F_\fp)\subseteq \F_q[\iota(M)]$. 
		\item[(iii)] There exists an eigenvalue $\la$ of $M$ (in $\oF_q$) such that $\F_q(\la)\supseteq \F_\fp$. 
	\end{itemize}
	Equivalently, (ii) fails if and only if every eigenvalue of $M$ lies in a proper subfield of $\F_\fp$ over $\F_q$. 
\end{prop}
\begin{proof}
	Let $m(x)\in \F_q[x]$ denote the minimal polynomial of $\iota(M)$ over $\F_q$, and let $m(x)=p_1(x)\cdots p_s(x)$ 
	be its factorization into distinct irreducible factors over $\F_q$, with $d_i=\deg(p_i)$. Then 
	\[
	\F_q[\iota(M)]\cong \F_q[x]/(m(x)) \cong \prod_{i=1}^s \F_{q^{d_i}}. 
	\]
	Since $\F_\fp\cong \F_{q^d}$, an embedding $\F_\fp\hookrightarrow \F_q[\iota(M)]$ exists if and only if $\F_{q^d}$ 
	embeds into some factor $\F_{q^{d_i}}$, which occurs if and only if $d\mid d_i$ for some $i$.  Thus, condition (ii) holds if and only if 
	some eigenvalue $\la$ of $\iota(M)$ satisfies $\F_\fp\subseteq \F_q(\la)$. 
	
	Finally, the eigenvalues of $\iota(M)$ are exactly the $\F_q$-conjugates of the eigenvalues of $M$. Since $\F_q$-conjugates 
	have the same minimal polynomial over $\F_q$, condition (iii) holds for an eigenvalue of $\iota(M)$ if and only if it holds 
	for an eigenvalue of $M$. 
\end{proof}

The same argument yields the following generalization.

\begin{prop}\label{propCM3}
	Let $M_1, \dots, M_k\in \Mat_r(\F_\fp)$. Then 
	\[
	C_{\Mat_{rd}(\F_q)}(\iota(M_1), \dots, \iota(M_k)) = \iota\left(C_{\Mat_r(\F_\fp)}(M_1, \dots, M_k)\right)
	\]
	if and only if $\iota(\F_\fp)\subseteq \F_q[\iota(M_1),\dots, \iota(M_k)]$. 
\end{prop}

Now returning to the message space $\cM$ in \eqref{eqcM}, we note that $\pi=\tau^n\in \cM$ once $rd>\ell$.  
We are assuming that $m_\phi(x)$ is irreducible modulo $\fp$, so the eigenvalues of $\pi$ acting on $\phi[\fp]$ 
are not in $\F_\fp$. Applying Propositions \ref{propCM} and \ref{propCM2}, we conclude that $\iota(\F_\fp)\subseteq \F_q[\iota(\pi)]$. 
Thus, by Proposition \ref{propCM3}, 
\[
C_{\End_{\F_\fp}(\phi[\fp])}(\cM) = C_{\End_{\F_q}(\phi[\fp])}(\cM).  
\]











%---------------------------------------------------------
%\renewcommand{\bibliofont}{\normalsize}
\bibliographystyle{amsalpha}
\bibliography{Bibliography.bib}

\end{document}

